\documentclass[11pt,a4paper]{article}

\usepackage[margin=0.82in]{geometry}
\usepackage{fontspec}
\usepackage{xeCJK}
\xeCJKsetup{CJKspace=true}
\setmainfont{DejaVu Serif}
\setsansfont{DejaVu Sans}
\setmonofont{DejaVu Sans Mono}
\setCJKmainfont{Noto Sans CJK KR}
\setCJKsansfont{Noto Sans CJK KR}
\setCJKmonofont{Noto Sans Mono CJK KR}
\usepackage{amsmath,amssymb,mathtools,bm}
\usepackage{booktabs}
\usepackage{longtable}
\usepackage{tabularx}
\usepackage{array}
\usepackage{enumitem}
\usepackage{xcolor}
\usepackage{hyperref}
\usepackage{listings}
\usepackage{titlesec}
\usepackage{fancyhdr}
\usepackage{microtype}
\usepackage{ragged2e}
\usepackage{graphicx}
\usepackage{float}
\usepackage{setspace}
\usepackage{multirow}

\hypersetup{
  colorlinks=true,
  linkcolor=blue!55!black,
  urlcolor=blue!55!black,
  citecolor=blue!55!black,
  pdftitle={Topology-Distilled, Error-Triggered Global-Local Wind Dynamics for Static 3D Gaussian Thin Surfaces},
  pdfauthor={Wind3DGS idea sketch}
}

\definecolor{notegray}{RGB}{245,246,248}
\definecolor{darkblue}{RGB}{30,58,138}
\definecolor{darkgreen}{RGB}{22,101,52}
\definecolor{darkred}{RGB}{153,27,27}
\definecolor{darkorange}{RGB}{154,52,18}
\definecolor{darkpurple}{RGB}{107,33,168}
\definecolor{softblue}{RGB}{239,246,255}
\definecolor{softgreen}{RGB}{240,253,244}
\definecolor{softred}{RGB}{254,242,242}

\lstset{
  basicstyle=\ttfamily\small,
  breaklines=true,
  columns=fullflexible,
  frame=single,
  backgroundcolor=\color{notegray},
  xleftmargin=0.5em,
  xrightmargin=0.5em,
  framerule=0pt
}

\setlength{\headheight}{14pt}
\setlength{\tabcolsep}{4pt}
\emergencystretch=3em
\allowdisplaybreaks
\pagestyle{fancy}
\fancyhf{}
\lhead{Topology-Distilled Error-Triggered Wind Dynamics}
\rhead{2026-08-13}
\cfoot{\thepage}

\titleformat{\section}{\Large\bfseries\color{darkblue}}{\thesection}{0.8em}{}
\titleformat{\subsection}{\large\bfseries}{\thesubsection}{0.8em}{}
\titleformat{\subsubsection}{\normalsize\bfseries}{\thesubsubsection}{0.8em}{}

\makeatletter
\renewcommand*\l@subsection{\@dottedtocline{2}{1.5em}{3.3em}}
\renewcommand*\l@subsubsection{\@dottedtocline{3}{4.8em}{4.2em}}
\makeatother

\newcommand{\R}{\mathbb{R}}
\newcommand{\eps}{\varepsilon}
\newcommand{\diag}{\operatorname{diag}}
\newcommand{\orth}{\operatorname{orth}}
\newcommand{\cof}{\operatorname{cof}}
\newcommand{\clamp}{\operatorname{clamp}}
\newcommand{\sigmoid}{\operatorname{sigmoid}}
\newcommand{\argmin}{\operatorname*{arg\,min}}
\newcommand{\ours}{\textcolor{darkgreen}{\textbf{[제안]}}}
\newcommand{\physics}{\textcolor{darkblue}{\textbf{[물리식]}}}
\newcommand{\learned}{\textcolor{darkpurple}{\textbf{[학습]}}}
\newcommand{\deterministic}{\textcolor{darkgreen}{\textbf{[고정 연산]}}}
\newcommand{\constrained}{\textcolor{darkorange}{\textbf{[제약 풀이]}}}
\newcommand{\trainonly}{\textcolor{darkorange}{\textbf{[학습 전용]}}}
\newcommand{\runtime}{\textcolor{darkgreen}{\textbf{[Runtime]}}}
\newcommand{\optional}{\textcolor{darkblue}{\textbf{[선택]}}}
\newcommand{\risk}{\textcolor{darkred}{\textbf{[위험]}}}
\newcommand{\killgate}{\textcolor{darkred}{\textbf{[중단 기준]}}}
\newcommand{\baseline}{\textcolor{darkpurple}{\textbf{[비교]}}}
\newcolumntype{Y}{>{\RaggedRight\arraybackslash}X}
\newcolumntype{L}[1]{>{\RaggedRight\arraybackslash}p{#1}}

\title{\textbf{Topology-Distilled, Error-Triggered Global--Local Wind Dynamics}\\
\textbf{for Static 3D Gaussian Thin Surfaces}\\
\large 독립 아이디어 스케치: 물리식 기반 Global predictor와 선택적 Local missing-force corrector}
\author{Wind3DGS 연구 설계 문서}
\date{2026년 8월 13일}

\begin{document}
\maketitle

\begin{center}
\fbox{\parbox{0.95\linewidth}{
\textbf{한 문장 주장.}
독립적으로 재구성된 static 3DGS thin-surface asset에서 training-only mesh teacher로
material topology와 축약 물리 연산자를 distill하고,
runtime에는 현재 변형 표면의 공기력을 object-level reduced dynamics로 적분한 뒤,
global model의 미래 오차가 큰 영역에서만 global modal complement의 local missing-force correction을 계산한다.\\[0.45em]
\textbf{Global 원리.}
현재 위치, 속도, 법선, 면적과 spatial wind에서 quasi-steady aerodynamic force를 계산하고,
force/moment-preserving hyper-reduction으로 object-wide generalized force를 만든 뒤,
항상 켜진 mass--stiffness--damping system을 적분한다.\\[0.45em]
\textbf{Local 원리.}
학습된 gate는 local correction의 향후 비용 대비 이득을 예측한다.
선택된 patch에서만 network가 base model이 놓친 local force를 예측하고,
global mode와 겹치지 않는 작은 local physical system이 그 힘을 적분한다.
겹치는 patch는 force와 torque를 보존하도록 먼저 조립하며,
수정된 형상의 aerodynamic load 차이는 다시 Global로 되돌린다.\\[0.45em]
\textbf{Runtime 비사용 항목.}
Target mesh, full MPM/FEM, Navier--Stokes/CFD, frame-wise Gaussian 재학습은 사용하지 않는다.
RSH는 mainline에서 제거하고 별도의 공력 표현 비교 항목으로만 남긴다.\\[0.45em]
\textbf{초기 범위.}
고정 또는 compliant attachment를 가진 single-layer thin surfaces:
flag, pennant, hanging cloth, curtain strip, paper sheet, leaf, petal.
Self-contact, tearing, detached flight, full tree, two-way fluid field update는 초기 논문 범위가 아니다.
}}
\end{center}

\tableofcontents
\newpage

\section{문서 목적과 오늘 확정한 결정}

이 문서는 기존 RSH 중심 문서의 revision이 아니다.
2026년 8월 13일 논의에서 새로 정리한 연구 방향을 독립적인 문제 정의, 방법, 학습 프로토콜,
비교 실험, 선행연구 경계로 처음부터 다시 작성한 아이디어 스케치다.
기존 파일은 수정하지 않으며, 본 문서는 다음 다섯 결정을 출발점으로 삼는다.

\begin{enumerate}[leftmargin=1.7em]
  \item \textbf{Global motion은 축약 물리식으로 푼다.}
  물체 전체의 관성, 복원, 감쇠, attachment, 큰 굽힘과 저주파 진동은 항상 켜진 reduced structural dynamics가 담당한다.
  \item \textbf{Local motion도 물리식으로 적분한다.}
  Network는 위치를 직접 출력하지 않고, analytic/global/local base가 놓친 missing force만 예측한다.
  작은 local mass--stiffness--damping system이 그 힘을 시간 적분하여 displacement와 velocity를 만든다.
  \item \textbf{Gate는 deformation magnitude가 아니라 local correction의 실제 이득을 예측한다.}
  학습 정답은 global-only와 oracle-local counterfactual rollout의 향후 오차 차이로 만든다.
  \item \textbf{중복 제거는 두 채널로 분리한다.}
  Force label에서는 base가 이미 설명한 항을 빼고,
  motion space에서는 local basis를 global modal complement로 제한한다.
  실제 local reaction과 corrected-geometry aerodynamic delta는 제거하지 않고 Global에 feedback한다.
  \item \textbf{Patch를 각각 적분한 뒤 평균하지 않는다.}
  Overlap된 patch force를 먼저 하나의 보존적인 local force로 조립하고,
  그 뒤 global complement에 정의된 단일 local state를 적분한다.
\end{enumerate}

\subsection{핵심 연구 질문}

\begin{enumerate}[leftmargin=1.7em]
  \item Static 3DGS만 주어진 target에서 runtime mesh 없이 thin-surface material connectivity,
  attachment, area/mass, bending/stretching operator와 reduced bases를 복원할 수 있는가?
  \item Current-deformed Gaussian surface의 spatial aerodynamic load를 소수 sample만으로 평가하면서
  object-level generalized force와 net wrench를 보존할 수 있는가?
  \item Global reduced physics가 놓치는 local flutter와 fold를,
  필요한 patch에서만 missing-force network와 local physics로 복원할 수 있는가?
  \item Error-triggered Local이 always-on Local과 유사한 품질을 유지하면서
  실제 GPU kernel 실행과 wall-clock cost를 의미 있게 줄이는가?
  \item Local correction이 단순 visual displacement가 아니라
  corrected aerodynamic geometry와 structural coupling을 통해 Global state에 기계적으로 영향을 주는가?
\end{enumerate}

\section{문제 정의}

\subsection{Runtime 입력}

Target runtime의 입력은 다음이다.

\begin{equation}
  \mathcal I_{\mathrm{run}}
  =
  \left(
    G^0,\;
    W(\mathbf x,t),\;
    \vartheta,\;
    \mathcal C,\;
    s_{\mathrm{metric}}
  \right).
\end{equation}

\section{범위, 가정, 비주장}

\subsection{초기 물리 범위}

\begin{itemize}[leftmargin=1.5em]
  \item Single-layer 또는 locally single-layer로 볼 수 있는 thin surfaces.
  \item Hard 또는 compliant attachment가 있는 object.
  \item Moderate deformation과 flutter; 심한 tearing과 topology change 제외.
  \item Prescribed one-way wind: object deformation은 aerodynamic load를 바꾸지만,
  object가 surrounding wind field 자체를 다시 푸는 full two-way CFD/FSI는 제외.
  \item Contact와 self-contact는 첫 논문에서 제외하거나 teacher-only stress test로 제한.
\end{itemize}

\subsection{명시적으로 주장하지 않는 것}

\begin{itemize}[leftmargin=1.5em]
  \item 최초의 modal wind simulation.
  \item 최초의 Gaussian surface aerodynamic force.
  \item 최초의 mesh-free reduced 3DGS dynamics.
  \item 최초의 hierarchical/global--local Gaussian simulator.
  \item 최초의 learned residual physics.
  \item 최초의 adaptive local/full-space simulation.
  \item Full CFD 또는 engineering-grade aeroelastic prediction.
  \item Appearance SH 또는 affine Gaussian transport 자체의 신규성.
\end{itemize}

\section{핵심 상태와 물리량 소유권}

\subsection{Global--Local 상태 분해}

Sparse material anchors의 rest position을 \(x^0\in\R^{3K}\)라 한다.
Runtime deformation은 다음처럼 분해한다.

\begin{equation}
  u^t
  =
  \Phi q^t
  +
  \Psi_{\mathcal S_t}z^t,
  \qquad
  x^t=x^0+u^t,
  \qquad
  v^t=\Phi\dot q^t+\Psi_{\mathcal S_t}\dot z^t.
  \label{eq:state-decomposition}
\end{equation}

\(\Phi\in\R^{3K\times r_g}\)는 object-wide low-frequency basis,
\(\Psi_{\mathcal S_t}\)는 active 또는 decay 중인 patch basis를
overlap-aware하게 조립한 Local basis다.
Local basis는 다음 조건을 만족하도록 미리 구성한다.

\begin{equation}
  \Phi^\top M\Psi=0.
  \label{eq:mass-orthogonality}
\end{equation}

식~\eqref{eq:mass-orthogonality}의 의미는 Local이 Global이 이미 표현한
mass-weighted motion을 다시 만들지 않는다는 것이다.
가능하면 runtime 후처리로 반복 투영하기보다 offline basis construction에서 이 조건을 만족시킨다.

\subsection{물리량 owner 계약}

\begin{longtable}{L{0.22\linewidth}L{0.27\linewidth}L{0.43\linewidth}}
\toprule
물리량 & 유일한 owner & 계약 \\
\midrule
\endfirsthead
\toprule
물리량 & 유일한 owner & 계약 \\
\midrule
\endhead
Mass/inertia & Global/Local reduced physical system & Network가 inertia를 다시 예측하지 않는다. \\
\addlinespace
Linear 또는 corotational restoring & \(K_g,K_\ell,K_{g\ell}\) & Gate와 무관하게 존재한다. Dormant local coordinate는 0이고, 이미 활성화된 local energy는 물리적으로 감쇠한다. \\
\addlinespace
Damping & \(C_g,C_\ell,C_{g\ell}\) & Network residual과 중복 학습하지 않는다. \\
\addlinespace
Gravity & explicit external force & Teacher residual에서 반드시 제거한다. \\
\addlinespace
Attachment / support & constraint 또는 compliant force 중 하나 & 같은 attachment를 force와 projection으로 동시에 적용하지 않는다. \\
\addlinespace
Quasi-steady aero & current-surface analytic law & Relative wind, current normal/area로 매 frame 계산한다. \\
\addlinespace
Missing local force & learned residual network & Base analytic/structural force를 제거한 teacher defect만 학습한다. \\
\addlinespace
Global nonlinear correction & optional learned generalized-force residual & Mainline 필수 아님. Direct state prediction은 baseline이다. \\
\addlinespace
Overlap reconciliation & constrained force assembly & Network가 patch seam을 임의로 숨기게 두지 않는다. \\
\addlinespace
Gaussian transport & final total anchor deformation & Global과 Local이 Gaussian을 각각 두 번 움직이지 않는다. \\
\bottomrule
\end{longtable}

\subsection{Gate가 끄는 것과 끄지 않는 것}

Gate는 Global restoring, Global damping, gravity, attachment를 절대 끄지 않는다.
Gate가 제어하는 것은 다음이다.

\begin{itemize}[leftmargin=1.5em]
  \item dormant patch의 Local enrichment 시작 여부,
  \item active patch의 expensive missing-force network 실행 여부,
  \item active aerodynamic correction sample 평가 여부.
\end{itemize}

한 번 local energy를 가진 patch는 gate가 off가 되어도 즉시 \(z,\dot z\)를 0으로 reset하거나 freeze하지 않는다.
Network 입력만 끄고 작은 local physical system으로 decay를 계속한다.
Energy가 충분히 작고 예상 benefit도 낮을 때만 dormant state로 반환한다.

\section{Contribution hierarchy와 신규성 경계}

\subsection{권장 contribution}

\begin{enumerate}[leftmargin=1.7em]
  \item \textbf{Topology-distilled target-mesh-free Gaussian thin-shell ROM.}
  Unordered static 3DGS에서 oriented material adjacency, area/mass, attachment,
  bending/stretching operator와 Global/Local bases를 추론한다.
  \item \textbf{Current-state spatial aerodynamic hyper-reduction.}
  현재 변형된 surface normal, area, relative velocity, spatial wind를 사용하면서
  소수 sample로 generalized force와 net wrench를 보존한다.
  \item \textbf{Error-triggered complementary Local corrector.}
  Counterfactual future benefit로 active patch를 선택하고,
  learned missing force를 Global modal complement의 local physics로 적분하며,
  overlap 보존과 corrected-aero feedback을 결합한다.
\end{enumerate}

\subsection{개별 식과 exact chain의 구분}

\begin{longtable}{L{0.27\linewidth}L{0.18\linewidth}L{0.47\linewidth}}
\toprule
연산 & 개별 신규성 & 정확한 경계 \\
\midrule
\endfirsthead
\toprule
연산 & 개별 신규성 & 정확한 경계 \\
\midrule
\endhead
\(x=x^0+\Phi q\), reduced ODE & 낮음 & 고전 subspace dynamics와 동일 계열이다 \cite{barbic2005stvk}. \\
\addlinespace
Wind force를 modal force로 projection & 낮음 & Wind Projection Basis가 직접적인 선행이다 \cite{diener2009windbasis}. \\
\addlinespace
Current Gaussian normal/area aero & 낮음 & Gaussian Swaying이 직접적인 선행이다 \cite{yan2026gaussianswaying}. \\
\addlinespace
Static GS reduced deformation & 낮음--중간 & DynamicTree, FreeForm, PhysSkin과 가깝다 \cite{li2026dynamictree,xiang2026freeform,lei2026physskin}. \\
\addlinespace
Physics + learned residual & 낮음 & PGRD와 Learned Contact Corrections가 직접 경계를 만든다 \cite{patel2026pgrd,romero2021contact}. \\
\addlinespace
Adaptive local enrichment & 낮음 & Subspace Condensation과 Trading Spaces가 선행한다 \cite{teng2015condensation,trusty2024tradingspaces}. \\
\addlinespace
Global complement motion & 낮음 & Complementary dynamics 계열이 선행한다 \cite{benchekroun2023complementary}. \\
\addlinespace
제안 전체 chain & 조건부 중간 & Static GS thin surface, training-only topology distillation, current-state conservative aero ROM, learned benefit gate, sparse complementary missing-force, corrected-aero feedback의 정확한 연결은 조사 범위에서 발견하지 못했다. 단, 구성요소의 obvious composition 반론을 operator 설계와 실험으로 넘어야 한다. \\
\bottomrule
\end{longtable}

\subsection{Obvious composition을 넘기 위한 조건}

Gaussian Swaying의 force law를 FreeForm 또는 DynamicTree의 reduced solver에 연결하고
일반 residual network를 더하는 것만으로는 강한 novelty가 되기 어렵다.
다음 네 항목 중 적어도 세 항목이 실제 새 연산으로 구현되고,
각각의 ablation이 end-to-end 차이를 보여야 한다.

\begin{enumerate}[leftmargin=1.7em]
  \item 단순 Euclidean kNN이 아닌 oriented material topology와 thin-shell operator distillation.
  \item Current-deformed spatial aero를 net force/torque까지 보존하는 learned cubature.
  \item Counterfactual future benefit와 compute cost를 함께 예측하는 gate.
  \item Complement, overlap conservation, local structural coupling,
  corrected-aero delta를 포함하는 안정적인 two-way update.
\end{enumerate}

\section{세 단계 전체 파이프라인}

\begin{longtable}{L{0.17\linewidth}L{0.28\linewidth}L{0.47\linewidth}}
\toprule
Phase & 입력 & 출력 \\
\midrule
\endfirsthead
\toprule
Phase & 입력 & 출력 \\
\midrule
\endhead
Teacher generation & Source mesh, material/attachment, wind program, render cameras & High-resolution structural trajectories, optional fluid loads, multi-view images \\
\addlinespace
Distillation / training & Teacher trajectories, independently reconstructed static GS & Target-runtime object package, gate network, missing-force network \\
\addlinespace
Target runtime & Static GS, prescribed wind, user material/attachment controls & Global/Local state, deformed Gaussians, diagnostics and rendered frames \\
\bottomrule
\end{longtable}

\subsection{Target-runtime object package}

한 target asset의 setup 단계는 다음 package를 출력한다.

\begin{equation}
\begin{aligned}
  \mathcal P(G^0)=\{&
  x^0,\mathcal E,\mathcal H,\mathcal R,\;
  A^0,n^0,M,\;
  \vartheta,\mathcal C,\\
  &K,C,\Phi,\{B_r\},\Psi,\;
  \mathcal S_{\mathrm{aero}},\omega,\;
  \mathcal B_{\mathrm{GS}}
  \}.
\end{aligned}
\end{equation}

\(\mathcal E\)는 material adjacency,
\(\mathcal H\)는 bending/hinge-like 관계,
\(\mathcal R\)은 overlapping core--halo patches,
\(\mathcal S_{\mathrm{aero}}\)와 \(\omega\)는 aerodynamic cubature samples와 weights,
\(\mathcal B_{\mathrm{GS}}\)는 anchor-to-Gaussian binding이다.

\section{Teacher와 학습 데이터 생성}

\subsection{Teacher-S: 넓은 구조 상태를 위한 teacher}

주 데이터는 high-resolution nonlinear thin-shell FEM, cloth FEM,
또는 검증된 dense MPM에서 만든다.
Teacher-S는 다음 label을 제공한다.

\begin{itemize}[leftmargin=1.5em]
  \item mesh positions, velocities, accelerations,
  \item strain, curvature, bending/stretching energy,
  \item internal/external/reaction force breakdown,
  \item surface normal, area, deformation gradient,
  \item attachment reaction,
  \item optional Gaussian mean/covariance target.
\end{itemize}

Teacher-S가 사용하는 aerodynamic base가 runtime과 같은 quasi-steady law라면,
learned missing force는 basis truncation, reduced operator error,
nonlinear shell response와 topology/material distillation error만 학습한다고 주장한다.

\subsection{Teacher-FSI: 제한된 unsteady aerodynamic teacher}

Wake, vortex shedding, history-dependent flutter를 주장하려면
작은 asset subset에 CFD, LBM, 또는 실제 계측 기반 teacher가 추가로 필요하다.
이 경우에만 missing-force label에 다음 항이 들어갈 수 있다.

\begin{equation}
  \Delta f_{\mathrm{unsteady}}^*
  =
  f_{\mathrm{FSI/measurement}}^*
  -
  f_{\mathrm{quasi\mbox{-}steady}}^*.
\end{equation}

\risk Teacher-FSI 없이 \emph{wake를 학습한다} 또는
\emph{vortex shedding을 복원한다}고 주장하면 안 된다.
Teacher-S만 사용하는 MVP에서는 해당 항을 basis/structure correction으로 제한한다.

\subsection{Factorized data sweep}

\begin{equation}
  \mathcal D
  =
  \mathcal D_{\mathrm{shape}}
  \times\mathcal D_{\mathrm{material}}
  \times\mathcal D_{\mathrm{attachment}}
  \times\mathcal D_{\mathrm{wind}}
  \times\mathcal D_{\mathrm{GS}}
  \times\mathcal D_{\mathrm{initial}}.
\end{equation}

\begin{itemize}[leftmargin=1.5em]
  \item Shape: flag, strip, curtain, sheet, leaf, simple plant.
  \item Material: density, stretch/shear/bending stiffness, damping, anisotropy.
  \item Attachment: one edge, two points, partial compliant support, stem support.
  \item Wind direction/magnitude: yaw, elevation, weak-to-strong flutter range.
  \item Temporal wind: steady, sinusoidal gust, chirp, abrupt start/stop.
  \item Spatial wind: traveling gust, localized gust, nonuniform macro zones.
  \item GS quality: density, view count, covariance noise, opacity perturbation,
  holes, missing backside, floaters, reconstruction seed.
\end{itemize}

\subsection{Split protocol}

Frame-random split은 동일 trajectory의 이웃 frame을 train/test에 섞으므로 금지한다.
최소 split 단위는 trajectory이며, 주요 일반화 결과는 object-disjoint split으로 보고한다.

\begin{enumerate}[leftmargin=1.7em]
  \item Seen object / unseen wind.
  \item Seen object / unseen material.
  \item Seen object / unseen attachment.
  \item Unseen geometry within the same category.
  \item Unseen object category.
  \item Independent GS reconstruction of a seen source shape.
  \item Unseen object + material + wind combined OOD.
\end{enumerate}

\section{Topology와 reduced operator distillation}

\subsection{Surface-aware anchors}

Raw Gaussian count를 physical importance로 사용하지 않는다.
Opacity, covariance surface-likeness, multi-view contribution,
local isolation, boundary/curvature evidence를 함께 사용하여
surface-reliable Gaussian support를 만들고 \(K\ll N_G\)개의 material anchors를 선택한다.

각 anchor는 다음 attribute를 갖는다.

\begin{equation}
  a_k^0
  =
  \left(
  x_k^0,n_k^0,A_k^0,m_k,
  \vartheta_k,c_k,\eta_k
  \right),
\end{equation}

여기서 \(c_k\)는 attachment/compliance,
\(\eta_k\)는 reconstruction reliability다.
면적과 질량은 Gaussian별 중복 support를 그대로 합하지 않고
partition-of-unity로 총량을 보존한다.

\subsection{Oriented material topology}

Topology network는 Euclidean proximity만 출력하지 않는다.
후보 edge \((i,j)\)에 대해 다음을 예측한다.

\begin{equation}
  (\widehat y_{ij}^{\mathrm{mat}},
   \widehat \ell_{ij}^0,
   \widehat t_{ij},
   \widehat k_{ij}^{\mathrm{stretch}},
   \widehat k_{ij}^{\mathrm{shear}}).
\end{equation}

별도의 local relation은 bending/hinge-like coupling,
orientation consistency, disconnected-layer rejection을 나타낸다.
Training label은 mesh geodesic/material adjacency에서 오지만,
runtime graph는 target mesh를 요구하지 않는다.

\subsection{Mass, stiffness, damping}

추론된 graph와 physical fields에서 sparse matrices를 구성한다.

\begin{equation}
  M=\diag(m_1I_3,\ldots,m_KI_3),\qquad
  K=K_{\mathrm{stretch}}+K_{\mathrm{shear}}+K_{\mathrm{bend}}+K_{\mathrm{attach}},
\end{equation}

\begin{equation}
  C=\alpha_M M+\alpha_K K+C_{\mathrm{local}}.
\end{equation}

Positive mass, symmetric positive semidefinite stiffness/damping,
attachment null-space 제거는 construction으로 보장한다.
Network가 dense arbitrary \(M,C,K\)를 직접 출력하게 두지 않는다.

\subsection{Global basis}

Global basis는 generalized eigenproblem 또는 teacher POD와 physics mode의 hybrid에서 얻는다.

\begin{equation}
  K\phi_j=\lambda_jM\phi_j,
  \qquad
  \Phi=[\phi_1,\ldots,\phi_{r_g}],
  \qquad
  \Phi^\top M\Phi=I.
\end{equation}

Low-frequency bending, sway, twist와 attachment-driven motion을 우선한다.
Teacher POD를 사용할 때도 \(M\)-orthonormalization과 rest/attachment consistency를 적용한다.

\subsection{Patch Local basis와 complement}

Patch \(r\)의 raw local modes를 \(B_r\)라 하고,
\(R_r\)을 global anchor에서 patch anchor를 추출하는 operator,
\(W_r\)을 core--halo partition weight라 한다.
모든 patch basis를 먼저 global anchor space에 조립한다.

\begin{equation}
  B_{\mathrm{raw}}
  =
  [R_1^\top W_1B_1,\ldots,R_P^\top W_PB_P].
\end{equation}

Global span 제거 operator는

\begin{equation}
  P_\perp
  =
  I-\Phi(\Phi^\top M\Phi)^{-1}\Phi^\top M.
\end{equation}

최종 basis는

\begin{equation}
  \Psi=\orth_M(P_\perp B_{\mathrm{raw}}),
  \qquad
  \Phi^\top M\Psi=0.
\end{equation}

중요한 순서는
\(\text{patch basis 생성}\rightarrow\text{overlap 조립}\rightarrow\text{complement}\rightarrow\text{orthonormalize}\)다.
Patch별로 독립 projection과 integration을 하고 나중에 평균내지 않는다.

\subsection{Reduced block operators}

\begin{equation}
  M_{gg}=\Phi^\top M\Phi,
  \qquad
  M_{\ell\ell}=\Psi^\top M\Psi,
  \qquad
  M_{g\ell}=\Phi^\top M\Psi\approx0,
\end{equation}

\begin{equation}
  K_{gg}=\Phi^\top K\Phi,\quad
  K_{\ell\ell}=\Psi^\top K\Psi,\quad
  K_{g\ell}=\Phi^\top K\Psi,
\end{equation}

그리고 \(C\)에도 같은 block을 만든다.
\(K_{g\ell}\)와 \(C_{g\ell}\)가 0이 아니면 Local state의 structural reaction이 Global에 전달된다.
정확한 eigenmode 분할로 cross block이 0이면 structural reaction contribution은 주장하지 않고,
corrected-aero feedback만 유지한다.

\subsection{Gaussian binding}

각 Gaussian은 소수 anchor의 affine MLS support에 고정된다.

\begin{equation}
  \mathcal B_j
  =
  \{(k,w_{jk},\xi_{jk}^0)\}_{k\in\mathcal N_j}.
\end{equation}

Binding은 identity, translation, rigid rotation,
in-plane affine deformation을 tolerance 내에서 재현해야 한다.
Setup 후 frame마다 neighbor를 다시 고르지 않는다.

\section{Gate와 missing-force 학습 target}

\subsection{Teacher state projection}

Teacher anchor displacement \(u^{*,t}\)에서 Global coordinate를 얻는다.

\begin{equation}
  q^{*,t}
  =
  (\Phi^\top M\Phi)^{-1}\Phi^\top M u^{*,t}.
\end{equation}

Global reconstruction을 제거한 residual을 Local complement에 투영한다.

\begin{equation}
  z^{*,t}
  =
  (\Psi^\top M\Psi)^{-1}
  \Psi^\top M
  (u^{*,t}-\Phi q^{*,t}).
\end{equation}

\subsection{Missing-force label}

Known local physics가 teacher \(z^*\)를 만들기 위해 필요로 하는 force를 inverse dynamics로 계산한다.

\begin{equation}
\begin{aligned}
  \Delta Q_{\ell}^{*,t}
  =&\;
  M_{\ell\ell}\ddot z^{*,t}
  +C_{\ell\ell}\dot z^{*,t}
  +K_{\ell\ell}z^{*,t}\\
  &+C_{\ell g}\dot q^{*,t}
  +K_{\ell g}q^{*,t}
  -\Psi^\top
  \left(
    f_{\mathrm{aero}}^{\mathrm{base},*,t}
    +f_{\mathrm{gravity}}^t
    +f_{\mathrm{other}}^t
  \right).
  \label{eq:missing-force-label}
\end{aligned}
\end{equation}

식~\eqref{eq:missing-force-label}은 teacher total acceleration을 그대로 residual로 쓰지 않는다.
Mass, stiffness, damping, gravity, analytic aero, explicit attachment가 이미 설명한 부분을 모두 한 번씩 제거한다.
Teacher가 force breakdown을 직접 기록하면 이를 우선 사용한다.
Position만 있다면 smoothed differentiation과 differentiable rollout loss를 같이 사용한다.

\subsection{Gate benefit label}

Gate label은 단순 strain threshold가 아니다.
Horizon \(H\) 동안 Global-only rollout과
patch \(r\)의 oracle Local correction을 켠 counterfactual rollout을 비교한다.

\begin{equation}
  b_r^t
  =
  \frac{
    [\mathcal E_H^{G}(r,t)
    -\mathcal E_H^{G+L_r}(r,t)]_+
  }{
    C_r^{L}+\eps
  }.
  \label{eq:benefit-label}
\end{equation}

\(\mathcal E_H\)는 future position, velocity, normal, curvature,
strain, spectrum error의 weighted sum이고,
\(C_r^L\)은 실제 또는 예측된 Local 실행 비용이다.
식~\eqref{eq:benefit-label}은 \emph{Local이 필요해 보이는가}가 아니라
\emph{Local을 켜면 비용 대비 미래 오차가 실제로 얼마나 줄어드는가}를 정답으로 만든다.

모든 patch counterfactual이 너무 비싸면 다음 근사를 사용한다.

\begin{enumerate}[leftmargin=1.7em]
  \item Teacher complement energy가 큰 후보 patch만 선별한다.
  \item 강하게 겹치는 patch를 group으로 묶는다.
  \item Short horizon \(H\)로 benefit을 계산하고,
  일부 sequence에서 long-horizon calibration을 수행한다.
\end{enumerate}

\subsection{Runtime input leakage 방지}

Gate input은 runtime에서 관측 가능한 값만 사용한다.

\begin{equation}
  \chi_r^t
  =
  \left(
  q^t,\dot q^t,\;
  z_{\mathcal N(r)}^t,\dot z_{\mathcal N(r)}^t,\;
  \epsilon_r^t,\kappa_r^t,\;
  W_r^t,\;
  \vartheta_r,\mathcal C_r,\;
  h_r^t
  \right).
\end{equation}

Teacher future state, teacher residual, teacher force를 gate input으로 넣지 않는다.
그 값들은 label 생성에만 사용한다.

\section{Runtime 14단계 개요}

\begin{longtable}{L{0.055\linewidth}L{0.20\linewidth}L{0.18\linewidth}L{0.24\linewidth}L{0.24\linewidth}}
\toprule
\# & 단계 & Solver 종류 & 입력 & 출력 \\
\midrule
\endfirsthead
\toprule
\# & 단계 & Solver 종류 & 입력 & 출력 \\
\midrule
\endhead
1 & 이전 상태 읽기 & \deterministic & \(q,\dot q,z,\dot z\), gate memory & 현재 frame 초기 상태 \\
2 & 현재 표면 복원 & \deterministic & State, bases, binding & Position, velocity, normal, area \\
3 & 상대풍/공기력 & \physics & Surface state, prescribed wind & Sample aerodynamic force \\
4 & Global force 축약 & \constrained & Sample force, cubature, \(\Phi\) & Generalized load \(Q_g\) \\
5 & Global predictor & \physics & \(Q_g,M_g,C_g,K_g\) & \(\tilde q,\dot{\tilde q}\) \\
6 & Local 필요성 판단 & \learned & Global predictor, patch features/history & Benefit, uncertainty, active set \\
7 & Local force proposal & \physics+\learned & Active patch state, wind, material & Analytic base aero + two missing-force channels \\
8 & Local physical update & \physics & Assembled local force, \(M_\ell,C_\ell,K_\ell\) & \(z,\dot z\) \\
9 & Global 중복 제거 & \deterministic & Patch basis/force, \(\Phi,M\) & Complement-only local field \\
10 & Overlap 보존 조립 & \constrained & Three-channel patch proposals, overlap graph & Channel-preserving conservative force \\
11 & Global feedback/corrector & \physics & Cross reaction, corrected aero delta & Corrected \(q,\dot q\) \\
12 & 최종 anchor state & \deterministic & Corrected \(q,z\), bases & Anchor state/deformation gradient \\
13 & Gaussian transport/render & \deterministic & Anchor state, GS binding & \(G^t\), rendered image \\
14 & 상태/hysteresis 갱신 & \deterministic & Energy, benefit, uncertainty & Next state, dormant/active/decay set \\
\bottomrule
\end{longtable}

\paragraph{순서에 대한 구현 주의.}
위 표는 이해를 위한 논리 순서다.
실제 구현에서는 7번 patch force proposal을 10번에서 먼저 조립하고 9번 complement를 적용한 뒤,
8번의 단일 local physical update를 수행한다.
즉 \emph{patch별 적분 후 overlap 평균}은 금지한다.
상세 module은 실제 계산 순서와 데이터 의존성을 명시한다.

\section{Module 1: 이전 Global/Local 상태 읽기}

\subsection*{푸는 질문}

\emph{현재 바람만 보고 새 위치를 만드는 것이 아니라,
이전 프레임의 관성, 진동 phase, local 잔류 에너지를 어떻게 이어갈 것인가?}

\subsection*{Solver 분류}

\deterministic State gather와 buffer validation이다. Network와 새로운 물리 solve는 없다.

\subsection*{입력}

\begin{equation}
  s^{t-1}
  =
  \left(
  q^{t-1},\dot q^{t-1},
  z^{t-1},\dot z^{t-1},
  \mathcal A_{t-1},
  \mathcal D_{t-1},
  h^{t-1}
  \right).
\end{equation}

\(\mathcal A\)는 missing-force network를 실행하는 active patch,
\(\mathcal D\)는 network는 꺼졌지만 local energy를 물리적으로 감쇠시키는 decay patch,
\(h\)는 gate hysteresis와 선택적인 aerodynamic memory다.

\subsection*{적용 연산}

\begin{itemize}[leftmargin=1.5em]
  \item State dimension, active basis index, timestep, asset package version을 검사한다.
  \item 새 control event가 있으면 attachment/material-dependent reduced operator를 갱신한다.
  \item Wind 변화만으로 topology, bases, Gaussian binding을 다시 만들지 않는다.
  \item Hard attachment가 이동했다면 prescribed support position/velocity를 읽는다.
\end{itemize}

\subsection*{출력과 다음 단계}

검증된 \(s^{t-1}\)가 Module 2의 current surface reconstruction으로 넘어간다.
State를 저장하기 때문에 resonance, wind-stop recovery, damping envelope가 프레임 사이에 유지된다.

\subsection*{필수 검증}

\begin{itemize}[leftmargin=1.5em]
  \item Zero-input state copy가 bitwise 또는 tolerance 내에서 일관적인가?
  \item Patch가 active에서 decay로 바뀌어도 \(z,\dot z\)가 reset되지 않는가?
  \item Asset/control 변경 시 오래된 operator와 새 state가 혼용되지 않는가?
\end{itemize}

\section{Module 2: 현재 anchor와 surface state 복원}

\subsection*{푸는 질문}

\emph{작은 Global/Local coordinate를 실제 공간의 위치, 속도, 법선, 면적으로 어떻게 바꿀 것인가?}

\subsection*{Solver 분류}

\deterministic Modal lifting, affine MLS, deformation-gradient geometry update다.

\subsection*{입력}

\begin{itemize}[leftmargin=1.5em]
  \item \(q^{t-1},\dot q^{t-1},z^{t-1},\dot z^{t-1}\),
  \item rest anchors \(x^0\), Global basis \(\Phi\), current Local basis \(\Psi_{\mathcal S}\),
  \item aerodynamic samples와 anchor support,
  \item rest normal \(n_s^0\), rest area \(A_s^0\).
\end{itemize}

\subsection*{위치와 속도}

\begin{equation}
  x^{t-1}=x^0+\Phi q^{t-1}+\Psi_{\mathcal S}z^{t-1},
\end{equation}

\begin{equation}
  v^{t-1}=\Phi\dot q^{t-1}+\Psi_{\mathcal S}\dot z^{t-1}.
\end{equation}

수식의 의미는 Global 큰 움직임과 Local detail을 한 번만 합쳐
현재 anchor 위치와 표면 속도를 만든다는 것이다.

\subsection*{Deformation gradient와 area-normal}

Sample \(s\) 주변 anchor의 rest/current correspondence에서
rank-2 surface-aware affine map \(F_s\)를 least squares로 맞춘다.
Current area-normal vector는 cofactor relation을 사용한다.

\begin{equation}
  \widetilde n_s^t
  =
  \cof(F_s^t) A_s^0 n_s^0,
  \qquad
  A_s^t=\|\widetilde n_s^t\|_2,
  \qquad
  n_s^t=\frac{\widetilde n_s^t}{A_s^t+\eps}.
  \label{eq:area-normal}
\end{equation}

식~\eqref{eq:area-normal}은 normal과 area를 별도 heuristic으로 갱신하지 않고
동일한 deformation map에서 일관되게 얻는다.
Degenerate stretch에서는 singular value clamp와 orientation continuity를 적용한다.

\subsection*{출력과 다음 단계}

\begin{equation}
  \mathcal G_{\mathrm{surf}}^{t-1}
  =
  \{x_s^{t-1},v_s^{t-1},F_s^{t-1},n_s^{t-1},A_s^{t-1}\}_{s\in\mathcal S_{\mathrm{aero}}}.
\end{equation}

이 surface state가 Module 3의 relative-wind/aerodynamic force 계산으로 넘어간다.
동일 \(F_s\)는 Module 13의 Gaussian covariance transport에도 재사용한다.

\subsection*{필수 검증}

\begin{itemize}[leftmargin=1.5em]
  \item Identity와 rigid rotation에서 area가 보존되고 normal이 정확히 회전하는가?
  \item In-plane scale에서 area가 예상 비율로 변하는가?
  \item Aero normal과 rendering/transport normal이 동일한 convention을 쓰는가?
  \item Normal sign이 frame 사이에 갑자기 뒤집히지 않는가?
\end{itemize}

\section{Module 3: Relative wind와 quasi-steady surface force}

\subsection*{푸는 질문}

\emph{현재 위치에서 움직이고 있는 표면이 실제로 느끼는 바람과
그 바람이 만드는 기본 external force는 무엇인가?}

\subsection*{Solver 분류}

\physics Prescribed wind interpolation과 quasi-steady aerodynamic law다.
Runtime에 Navier--Stokes, LBM, full CFD를 풀지 않는다.

\subsection*{입력}

\begin{itemize}[leftmargin=1.5em]
  \item Module 2의 \(x_s,v_s,n_s,A_s\),
  \item prescribed wind field \(W(\mathbf x,t)\),
  \item density \(\rho_{\mathrm{air}}\),
  \item drag/friction/lift coefficients \(\theta_s^{\mathrm{aero}}\),
  \item optional exposure/reliability \(e_s\).
\end{itemize}

\subsection*{Relative wind}

\begin{equation}
  r_s^t
  =
  W(x_s^t,t)-v_s^t,
  \qquad
  \widehat r_s^t
  =
  \frac{r_s^t}{\|r_s^t\|_2+\eps}.
  \label{eq:relative-wind}
\end{equation}

표면이 바람과 같은 방향으로 움직이면 실제로 느끼는 속도가 줄고,
반대 방향으로 움직이면 증가한다.
Absolute wind만 쓰는 baseline은 이 효과를 제거하는 ablation이다.

\subsection*{Surface aerodynamic law}

구체적인 drag/lift convention은 teacher와 runtime에서 동일하게 고정한다.
문서 수준에서는 다음 generic calibrated map으로 쓴다.

\begin{equation}
  f_{a,s}^t
  =
  \frac12\rho_{\mathrm{air}}
  A_s^t
  \|r_s^t\|_2^2
  e_s^t
  \mathcal C_{\mathrm{aero}}
  \left(
    \widehat r_s^t,n_s^t,
    \theta_s^{\mathrm{aero}}
  \right).
  \label{eq:surface-aero}
\end{equation}

\(\mathcal C_{\mathrm{aero}}\)는 normal pressure/drag,
tangential friction, optional lift를 반환한다.
MVP에서는 two-sided thin-sheet law를 사용하고,
front/back 계수가 다르면 oriented normal label을 유지한다.

\subsection*{수식이 풀지 않는 것}

식~\eqref{eq:surface-aero}은 local quasi-steady force다.
다음 현상은 직접 풀지 않는다.

\begin{itemize}[leftmargin=1.5em]
  \item object가 wind field에 만드는 wake,
  \item vortex shedding의 fluid state,
  \item long aerodynamic memory,
  \item dense self-shadowing과 full pressure solve.
\end{itemize}

이 현상을 주장하려면 Teacher-FSI가 필요하고,
그 차이만 Module 7의 missing-force target에 포함한다.

\subsection*{출력과 다음 단계}

Sample force set \(\{f_{a,s}^t\}\)가 Module 4의 conservative hyper-reduction으로 넘어간다.
Active Local correction 이후에는 같은 식을 Module 11에서 active region에만 다시 평가한다.

\subsection*{필수 검증}

\begin{itemize}[leftmargin=1.5em]
  \item Zero relative wind에서 force가 정확히 0인가?
  \item Rigid co-rotation에서 force가 frame-equivariant한가?
  \item Surface velocity를 제거한 baseline과 강한 gust에서 차이가 나는가?
  \item Current normal/area 대신 rest 값을 쓴 baseline 대비 deformed-state force가 정확한가?
\end{itemize}

\section{Module 4: Conservative aerodynamic hyper-reduction}

\subsection*{푸는 질문}

\emph{모든 Gaussian/anchor의 공기력을 평가하지 않고도
Global dynamics가 받는 generalized force와 물체 전체 wrench를 정확히 만들 수 있는가?}

\subsection*{Solver 분류}

\trainonly Empirical cubature sample/weight optimization,
\runtime weighted force evaluation과 Galerkin projection이다.
Cubature 자체는 고전적이며 \cite{an2008cubature},
본 연구의 후보 기여는 current Gaussian thin surface와 spatial wind에 맞춘 wrench-preserving construction이다.

\subsection*{Full reference}

모든 surface elements를 평가한다면 Global generalized force는

\begin{equation}
  Q_{g,\mathrm{aero}}^t
  =
  \sum_{s=1}^{N_S}
  J_{s,g}^{t\top} f_{a,s}^t,
  \label{eq:full-generalized-aero}
\end{equation}

여기서 \(J_{s,g}\)는 Global coordinate 변화가 sample 위치에 만드는 velocity/displacement Jacobian이다.
Linear lifting이면 \(\Phi_s\)와 동일 계열이고,
corotational/current map에서는 현재 frame Jacobian을 사용한다.

\subsection*{Hyper-reduced evaluation}

\begin{equation}
  \widehat Q_{g,\mathrm{aero}}^t
  =
  \sum_{s\in\mathcal S_c}
  \omega_s
  J_{s,g}^{t\top} f_{a,s}^t,
  \qquad
  |\mathcal S_c|\ll N_S.
  \label{eq:hyper-aero}
\end{equation}

Sample set \(\mathcal S_c\)와 nonnegative weights \(\omega_s\)는
다양한 deformation/wind training snapshots에서 다음 목표를 함께 최소화한다.

\begin{equation}
\begin{aligned}
  \min_{\mathcal S_c,\omega\ge0}
  \sum_m
  &\left\|
  Q_{g,m}^{\mathrm{full}}
  -\widehat Q_{g,m}
  \right\|_{W_Q}^2\\
  +\lambda_F&
  \left\|
  F_m^{\mathrm{full}}
  -\widehat F_m
  \right\|_2^2
  +\lambda_\tau
  \left\|
  \tau_m^{\mathrm{full}}
  -\widehat\tau_m
  \right\|_2^2.
  \label{eq:cubature-fit}
\end{aligned}
\end{equation}

Net force와 torque는

\begin{equation}
  F=\sum_s f_{a,s},
  \qquad
  \tau=\sum_s(x_s-c)\times f_{a,s}
\end{equation}

로 계산한다.
Centroid \(c\) convention은 teacher, training, runtime에서 동일해야 한다.

\subsection*{입력}

Module 3의 current sample forces,
offline sample indices/weights,
current Global Jacobian, object center다.

\subsection*{출력과 다음 단계}

\begin{equation}
  \widehat Q_{g,\mathrm{aero}}^t,\quad
  \widehat F_{\mathrm{aero}}^t,\quad
  \widehat\tau_{\mathrm{aero}}^t,\quad
  \epsilon_{\mathrm{cub}}^t.
\end{equation}

Generalized force는 Module 5로,
estimated cubature error와 sample confidence는 Module 6 gate feature로 넘어간다.

\subsection*{왜 network가 아닌가}

한 frame에 현재 wind query 하나만 필요하므로
모든 방향 coefficient를 생성하는 state-conditioned RSH보다
현재 sample force를 직접 평가하고 합산하는 편이 자연스럽다.
RSH, directional MLP, LUT는 E2 공력 표현 비교에만 남긴다.

\subsection*{필수 검증}

\begin{itemize}[leftmargin=1.5em]
  \item Full samples 대비 generalized-force, net-force, net-torque error.
  \item Random/FPS/nonconservative learned sampling과의 비교.
  \item Sample count 대 actual GPU latency의 Pareto curve.
  \item Held-out deformation과 traveling spatial gust에서의 error.
\end{itemize}

\section{Module 5: Always-on Global reduced physics predictor}

\subsection*{푸는 질문}

\emph{현재 전체 힘을 받았을 때 물체의 큰 굽힘, sway, twist, 저주파 진동이
다음 시간에 어떻게 변하는가?}

\subsection*{Solver 분류}

\physics Reduced second-order structural dynamics와 implicit time integration이다.
기본형에는 direct next-state network를 쓰지 않는다.

\subsection*{입력}

\begin{itemize}[leftmargin=1.5em]
  \item 이전 \(q^{t-1},\dot q^{t-1}\),
  \item Module 4의 \(\widehat Q_{g,\mathrm{aero}}^t\),
  \item reduced \(M_{gg},C_{gg},K_{gg}\),
  \item gravity, attachment/handle, explicit external force,
  \item 이전 Local state가 만드는 cross-block term,
  \item optional global nonlinear force correction.
\end{itemize}

\subsection*{Global equation}

\begin{equation}
\begin{aligned}
  M_{gg}\ddot q
  +C_{gg}\dot q
  +K_{gg}q
  +g_{\mathrm{nl}}(q)
  =&\;
  Q_{g,\mathrm{aero}}
  +Q_{g,\mathrm{gravity}}
  +Q_{g,\mathrm{handle}}\\
  &-C_{g\ell}\dot z
  -K_{g\ell}z
  +\Delta Q_{g,\theta}.
  \label{eq:global-dynamics}
\end{aligned}
\end{equation}

\(g_{\mathrm{nl}}\)은 corotational 또는 precomputed nonlinear reduced internal force다.
\(\Delta Q_{g,\theta}\)는 선택적인 learned generalized-force correction이며
MVP에서는 0으로 두고 ablation으로 필요성을 검사한다.

\subsection*{Time integration}

Implicit Newmark, backward Euler, 또는 implicit midpoint 중 하나를 고정한다.
예를 들어 backward Euler는

\begin{equation}
  q^t=q^{t-1}+\Delta t\,\dot q^t,
  \qquad
  \dot q^t=\dot q^{t-1}+\Delta t\,\ddot q^t
\end{equation}

를 식~\eqref{eq:global-dynamics}과 함께 푼다.
Global dimension \(r_g\)가 작으므로 same-frame corrector를 한 번 추가해도 비용이 제한적이다.

\subsection*{출력과 다음 단계}

Local을 아직 새로 적용하지 않은

\begin{equation}
  \widetilde q^t,\quad
  \dot{\widetilde q}^t,\quad
  \widetilde x_G^t=x^0+\Phi\widetilde q^t+\Psi z^{t-1}
\end{equation}

를 Global predictor라 한다.
이 값과 patch features가 Module 6 gate로 넘어간다.

\subsection*{왜 Global을 physics로 두는가}

Global system은 이미 저차원이므로 direct network로 대체했을 때의 속도 이득이 제한적이다.
반면 physics는 다음을 명시적으로 유지한다.

\begin{itemize}[leftmargin=1.5em]
  \item wind-stop recovery,
  \item resonance와 phase,
  \item mass/stiffness/damping control,
  \item unseen wind/material에 대한 구조적 외삽,
  \item long rollout energy behavior.
\end{itemize}

Direct neural Global은 중요한 baseline이지만 default method가 아니다.

\subsection*{필수 검증}

\begin{itemize}[leftmargin=1.5em]
  \item Zero wind에서 rest로 수렴하는가?
  \item Mode frequency와 damping envelope가 teacher와 일치하는가?
  \item Direct MLP/GNN/GRU Global보다 OOD와 장시간 rollout이 안정적인가?
  \item 동일 Global rank에서 sparse PD/XPBD보다 실제로 유리한가?
\end{itemize}

\section{Module 6: Learned future-benefit gate}

\subsection*{푸는 질문}

\emph{Global-only 결과를 그대로 사용하면 어느 patch에서 향후 오차가 커지고,
Local 계산 비용을 지불할 가치가 있는가?}

\subsection*{Solver 분류}

\learned Lightweight topology-aware gate network다.
모든 patch에서 실행되지만 expensive Local expert보다 훨씬 작아야 한다.

\subsection*{입력}

\begin{itemize}[leftmargin=1.5em]
  \item Global predictor \(q,\dot q,\widetilde x_G,\widetilde v_G\),
  \item patch strain, curvature, deformation rate,
  \item current relative wind와 wind gradient,
  \item cubature/topology/material uncertainty,
  \item past Local state energy와 recent activation history,
  \item basis truncation indicators.
\end{itemize}

\subsection*{출력}

\begin{equation}
  (\widehat b_r^t,\widehat\sigma_r^t,\widehat c_r^t)
  =
  G_\varphi(\chi_r^t),
\end{equation}

여기서 \(\widehat b\)는 예상 future benefit,
\(\widehat\sigma\)는 uncertainty,
\(\widehat c\)는 predicted Local runtime cost다.
Risk-aware score는 예를 들어

\begin{equation}
  s_r^t
  =
  \widehat b_r^t
  -\lambda_\sigma\widehat\sigma_r^t.
\end{equation}

\subsection*{Active set와 budget}

\begin{equation}
  \mathcal A_t
  =
  \operatorname{TopBudget}
  \left(
    \{r:s_r^t>\tau_{\mathrm{on}}\},
    B_{\mathrm{local}}
  \right).
\end{equation}

기존 active patch는 \(s_r^t<\tau_{\mathrm{off}}\)가 일정 시간 지속될 때만 decay로 이동한다.
\(\tau_{\mathrm{off}}<\tau_{\mathrm{on}}\)으로 hysteresis를 둔다.

\subsection*{연속 activation}

새 patch는 local coordinate를 0에서 시작하고,
force injection coefficient를 몇 frame에 걸쳐 증가시킨다.

\begin{equation}
  \gamma_r^{t+1}
  =
  \clamp(\gamma_r^t+\Delta\gamma_{\mathrm{on}},0,1)
\end{equation}

비활성 전환에서는

\begin{equation}
  \gamma_r^{t+1}
  =
  \clamp(\gamma_r^t-\Delta\gamma_{\mathrm{off}},0,1)
\end{equation}

로 줄인다.
Network evaluation을 시작하는 threshold와 full blend threshold를 분리하여
network가 처음 실행되는 frame에 full-scale force가 갑자기 나타나지 않게 한다.

\subsection*{출력과 다음 단계}

\(\mathcal A_t\), decay set \(\mathcal D_t\), \(\gamma_r^t\),
uncertainty와 budget diagnostics가 Module 7로 넘어간다.
Inactive patch의 missing-force kernel은 실제로 launch하지 않는다.

\subsection*{필수 검증}

\begin{itemize}[leftmargin=1.5em]
  \item Oracle benefit에 대한 recall, AUPRC, calibration error.
  \item 같은 active ratio의 random/curvature/strain gate와 비교.
  \item Activation jerk, patch toggle count, hysteresis ablation.
  \item Kernel profiler로 inactive patch가 실제 skip되는지 확인.
  \item Always-on Local과의 quality--latency Pareto.
\end{itemize}

\section{Module 7: Active-patch Local force proposal}

\subsection*{푸는 질문}

\emph{Global physics, analytic aero, reduced Local base가 설명하지 못한
추가적인 local force는 무엇인가?}

\subsection*{Solver 분류}

\physics Active region의 current analytic Local aero 평가와
\learned Shared patch-local equivariant GNN/MLP, optional recurrent memory를 묶는 orchestration이다.
최종 displacement나 next position을 직접 출력하지 않는다.

\subsection*{입력}

Active patch \(r\)마다 다음을 canonical/co-rotated local frame에 넣는다.

\begin{equation}
\begin{aligned}
  \xi_r^t=\{&
  R_r\widetilde x_G^t,\;
  R_r\widetilde v_G^t,\;
  R_r\Psi z^{t-1},\;
  R_r\Psi\dot z^{t-1},\\
  &\epsilon_r^t,\kappa_r^t,\dot\epsilon_r^t,\;
  W_r^t,\nabla W_r^t,\;
  \vartheta_r,\mathcal C_r,\;
  h_r^{t-1}\}.
\end{aligned}
\end{equation}

\subsection*{출력}

먼저 Module 3과 같은 analytic law를 active/halo sample에 적용하여
Local base aerodynamic force proposal \(f_{r,\mathrm{base}}^t\)를 계산한다.
이는 network output이 아니며 \texttt{analytic aero}가 physical owner다.

\begin{equation}
  (\widehat{\Delta f}_{r,\mathrm{aero}}^t,
   \widehat{\Delta f}_{r,\mathrm{struct}}^t,
   h_r^t)
  =
  N_\theta(\xi_r^t).
  \label{eq:missing-force-net}
\end{equation}

Teacher force breakdown이 신뢰 가능하면 aerodynamic residual과 structural/model residual을 분리한다.
그렇지 않으면 single total defect head를 사용하되 주장 범위를 축소한다.

Stage 7 orchestration의 최종 \texttt{PatchForceBatch}는
\begin{equation}
  \left\{
    f_{r,\mathrm{base}}^t,
    \widehat{\Delta f}_{r,\mathrm{aero}}^t,
    \widehat{\Delta f}_{r,\mathrm{struct}}^t
  \right\}_{r\in\mathcal A_t}
\end{equation}
의 세 channel을 분리해 담는다.
첫 channel은 analytic base aero이고 뒤의 두 channel만 learned missing force다.
Module 10과 9는 이 channel들을 합쳐 owner를 잃지 않고 각각 조립·투영한다.

\begin{itemize}[leftmargin=1.5em]
  \item Aerodynamic head는 external missing force이며 nonzero net wrench를 가질 수 있다.
  \item Structural head는 internal correction이므로 free object 부분에서는 net internal force/torque가 0이어야 한다.
  \item Attachment-related correction은 explicit reaction channel로 분리한다.
\end{itemize}

\subsection*{왜 network가 필요한가}

Network는 알려진 inertia, elasticity, damping을 대체하지 않는다.
다음처럼 analytic하게 정확히 풀면 runtime 비용이 커지거나 target 정보가 부족한 항만 근사한다.

\begin{itemize}[leftmargin=1.5em]
  \item Global/Local basis truncation,
  \item large-deformation thin-shell nonlinearity,
  \item topology/material distillation error,
  \item Teacher-FSI가 있을 때의 unsteady aerodynamic defect.
\end{itemize}

정확한 topology와 계산 여유가 있다면 active patch에 fine FEM/shell solve를 수행할 수 있다.
이는 반드시 포함해야 할 physics-only Local baseline이다.

\subsection*{출력과 다음 단계}

Patch force proposals는 즉시 적분하지 않는다.
Module 10의 overlap-aware constrained assembly와 Module 9의 complement decomposition을 거친 뒤
Module 8의 단일 Local physical system에 들어간다.

\subsection*{필수 검증}

\begin{itemize}[leftmargin=1.5em]
  \item Force target error와 multi-step rollout error를 함께 평가한다.
  \item Direct displacement network 대비 wind-stop decay와 phase 안정성을 비교한다.
  \item Physics-only Local 대비 learned defect의 순수 이득을 측정한다.
  \item Teacher-FSI가 없을 때 unsteady-aero claim이 결과표에 섞이지 않게 한다.
\end{itemize}

\section{Module 8: Single Local complementary physics update}

\subsection*{푸는 질문}

\emph{조립·투영된 Local force channel을 받았을 때 Global이 표현하지 못하는 Local state가
관성, 복원, 감쇠를 따라 어떻게 움직이는가?}

\subsection*{Solver 분류}

\physics Active/decay Local basis 위의 reduced second-order dynamics다.

\subsection*{입력}

\begin{itemize}[leftmargin=1.5em]
  \item Module 10과 9를 channel별로 거친
  \(Q_{\ell,\mathrm{base}}^{\perp,t}\),
  \(\Delta Q_{\ell,\mathrm{aero}}^{\perp,t}\),
  \(\Delta Q_{\ell,\mathrm{struct}}^{\perp,t}\),
  \item active/decay basis \(\Psi_{\mathcal S_t}\),
  \item \(M_{\ell\ell},C_{\ell\ell},K_{\ell\ell}\)와 cross blocks,
  \item predictor Global state \(\widetilde q,\dot{\widetilde q}\),
  \item 이전 \(z,\dot z\).
\end{itemize}

\subsection*{Local equation}

\begin{equation}
\begin{aligned}
  M_{\ell\ell}\ddot z
  +C_{\ell\ell}\dot z
  +K_{\ell\ell}z
  =&\;
  \Gamma_t\left(
  Q_{\ell,\mathrm{base}}^{\perp}
  +\Delta Q_{\ell,\mathrm{aero}}^{\perp}
  +\Delta Q_{\ell,\mathrm{struct}}^{\perp}
  \right)
  -C_{\ell g}\dot{\widetilde q}
  -K_{\ell g}\widetilde q.
  \label{eq:local-dynamics}
\end{aligned}
\end{equation}

\(\Gamma_t\)는 active local coordinates의 smooth activation weights다.
중요하게도 \(\Gamma_t\)는 \(M_{\ell\ell},C_{\ell\ell},K_{\ell\ell}\)를 곱하지 않는다.
즉 gate가 off되어도 이미 존재하는 Local state의 inertia, restoring, damping은 계속 작동한다.

\(Q_{\ell,\mathrm{base}}^{\perp}\)는 current analytic aero를 overlap-aware하게 조립하고
Global force span을 제거한 Local generalized load다.
두 \(\Delta Q^\perp\)만 learned missing-force channel이다.
세 channel은 서로 다른 owner/lineage를 유지한 채 같은 Module 10--9 경로를 거치고,
Module 8에서만 합산된다.
Decay patch에서는 새 excitation과 learned defect를 서서히 줄이되,
기존 \(z,\dot z\)는 물리적으로 감쇠한다.

\subsection*{출력}

\begin{equation}
  z^t,\quad\dot z^t,\quad
  E_{\ell}^t
  =
  \frac12\dot z^{t\top}M_{\ell\ell}\dot z^t
  +
  \frac12 z^{t\top}K_{\ell\ell}z^t.
\end{equation}

Local state는 Module 11의 corrected geometry와 Module 14의 decay 판단으로 넘어간다.

\subsection*{왜 patch별 integrate-then-average가 아닌가}

각 patch가 독립된 mass를 갖고 적분한 뒤 위치를 평균하면
overlap anchor의 실제 mass, energy, reaction이 정의되지 않는다.
본 방법은 patch force를 먼저 하나의 global anchor force로 조립하고,
하나의 complement basis system에서 한 번만 적분한다.

\subsection*{필수 검증}

\begin{itemize}[leftmargin=1.5em]
  \item Gate-off/zero-wind에서 Local energy가 증가하지 않는가?
  \item Patch activation/deactivation에서 position, velocity, acceleration이 연속적인가?
  \item Local physics 없이 direct displacement만 더한 방법보다 장시간 안정적인가?
  \item Timestep 변화에 대한 민감도가 허용 범위인가?
\end{itemize}

\section{Module 9: Global span과 Local force의 중복 분리}

\subsection*{푸는 질문}

\emph{Patch별 Local force channel 중 Global이 이미 담당하는 성분과
진짜 Local complement 성분을 어떻게 구분할 것인가?}

\subsection*{Solver 분류}

\deterministic Force-space projection과 diagnostics다.

\subsection*{Force-space complement}

Global virtual work가 0인 force를 만들기 위한 projector는

\begin{equation}
  \Pi_f^\perp
  =
  I-M\Phi(\Phi^\top M\Phi)^{-1}\Phi^\top.
  \label{eq:force-complement}
\end{equation}

Channel 집합
\(\mathcal C_\ell=\{\mathrm{base\ aero},\mathrm{missing\ aero},\mathrm{missing\ structural}\}\)의
조립된 raw force \(\widetilde f_{\ell,c}\)마다

\begin{equation}
  f_{\ell,c}^\perp
  =
  \Pi_f^\perp\widetilde f_{\ell,c},
  \qquad
  \Phi^\top f_{\ell,c}^\perp=0,
  \quad c\in\mathcal C_\ell.
\end{equation}

Local generalized force는

\begin{equation}
  Q_{\ell,c}^\perp
  =
  \Psi_{\mathcal S_t}^\top f_{\ell,c}^\perp,
  \quad c\in\mathcal C_\ell
\end{equation}

로 만든다.

\subsection*{제거하지 말아야 할 힘}

다음 두 항은 중복이 아니므로 버리지 않는다.

\begin{enumerate}[leftmargin=1.7em]
  \item \(K_{g\ell}z+C_{g\ell}\dot z\)가 나타내는 실제 structural coupling.
  \item Local correction이 표면 normal/area/velocity를 바꾸어 새로 생긴 aerodynamic generalized-force delta.
\end{enumerate}

Network가 별도의 Global missing-force head를 갖는다면
\(\Phi^\top\widetilde f\)는 그 head의 label/출력으로 명시적으로 보낸다.
Mainline Local head에서는 Global leakage를 loss와 projection으로 억제한다.

\subsection*{출력과 다음 단계}

세 channel의 \(Q_{\ell,c}^\perp\)는 owner/lineage를 보존한 채 Module 8로,
Global-span leakage norm

\begin{equation}
  \epsilon_{\mathrm{leak}}
  =
  \max_{c\in\mathcal C_\ell}
  \frac{\|\Phi^\top\widetilde f_{\ell,c}\|_2}
  {\|\widetilde f_{\ell,c}\|_2+\eps}
\end{equation}

은 diagnostics와 training loss로 넘어간다.

\section{Module 10: Overlap-aware conservative patch assembly}

\subsection*{푸는 질문}

\emph{여러 active patch가 같은 anchor에서 서로 다른 force를 제안할 때
seam, double counting, 가짜 translation/rotation 없이 하나의 force를 만드는가?}

\subsection*{Solver 분류}

\constrained Sparse scatter와 작은 constrained least-squares/KKT solve다.

\subsection*{초기 조립}

각 channel \(c\in\mathcal C_\ell\)을 합치기 전에 따로 조립한다.

\begin{equation}
  \widetilde f_c^t
  =
  \sum_{r\in\mathcal A_t}
  R_r^\top W_r
  \gamma_r^t
  f_{r,c}^t,
  \qquad c\in\mathcal C_\ell.
  \label{eq:raw-overlap-assembly}
\end{equation}

Partition weights는 displacement continuity에 유용하지만
힘과 torque 보존을 자동으로 보장하지 않는다.
따라서 식~\eqref{eq:raw-overlap-assembly}만으로 끝내지 않는다.

\subsection*{Constrained correction}

\begin{equation}
  f_{\mathrm{asm},c}
  =
  \argmin_f
  \|f-\widetilde f_c\|_{W_f}^2
  \quad
  \text{s.t.}\quad
  A_{\mathrm{cons},c}f=b_{\mathrm{cons},c}.
  \label{eq:conservative-assembly}
\end{equation}

Constraint는 channel별로 다르다.

\begin{itemize}[leftmargin=1.5em]
  \item Structural/internal residual:
  free internal net force/torque 0, attachment reaction은 별도 기록.
  \item Aerodynamic/external residual:
  overlap 조립 전 patch proposal의 unique resultant force/torque를 유지.
  \item Shared anchor:
  동일 DOF에 대한 action--reaction과 velocity continuity.
\end{itemize}

가능하면 structural residual을 antisymmetric edge force로 parameterize하고,
최종 assembled field에서 다시 net force/torque를 검사한다.

\subsection*{실제 계산 순서}

\begin{enumerate}[leftmargin=1.7em]
  \item Module 7 patch proposals 생성.
  \item 본 Module 10에서 force-space overlap 조립.
  \item Module 9에서 Global-span 분해.
  \item Module 8에서 단일 Local system 적분.
\end{enumerate}

\subsection*{출력과 다음 단계}

Channel별 conservative \(f_{\mathrm{asm},c}\), reaction wrench,
overlap constraint residual이 Module 9와 Module 11로 넘어간다.

\subsection*{필수 검증}

\begin{itemize}[leftmargin=1.5em]
  \item 단순 averaging 대비 patch boundary displacement/velocity jump.
  \item 최종 assembled net force/torque error.
  \item Free object에서 artificial translation/rotation drift.
  \item Patch partition이 달라져도 결과가 크게 달라지지 않는가?
\end{itemize}

\section{Module 11: Local-to-Global aeroelastic feedback와 corrector}

\subsection*{푸는 질문}

\emph{Local flutter/fold가 표면 방향과 구조 반작용을 바꾸었을 때
그 변화가 물체 전체의 다음 움직임에 어떻게 영향을 주는가?}

\subsection*{Solver 분류}

\physics Active-region aerodynamic reevaluation과 작은 reduced Global corrector다.

\subsection*{Predictor와 corrected geometry}

\begin{equation}
  x_{\mathrm{pred}}^t
  =
  x^0+\Phi\widetilde q^t+\Psi z^{t-1},
\end{equation}

\begin{equation}
  x_{\mathrm{corr}}^t
  =
  x^0+\Phi\widetilde q^t+\Psi z^t.
\end{equation}

Active/halo region에서 Module 2와 3의 normal/area/relative-wind force를 다시 평가한다.

\subsection*{Aerodynamic delta}

\begin{equation}
  \Delta Q_g^{\mathrm{aero},t}
  =
  \Phi^\top
  \left[
    f_a(x_{\mathrm{corr}}^t,v_{\mathrm{corr}}^t,W^t)
    -
    f_a(x_{\mathrm{pred}}^t,v_{\mathrm{pred}}^t,W^t)
  \right].
  \label{eq:corrected-aero-delta}
\end{equation}

Same-frame corrector에는 식~\eqref{eq:corrected-aero-delta}의 차이만 추가한다.
Module 5에서 이미 사용한 base aerodynamic force 전체를 다시 넣으면 double counting이다.

\subsection*{Structural coupling}

\begin{equation}
  Q_{g\leftarrow\ell}^{\mathrm{current},t}
  =
  -C_{g\ell}\dot z^t
  -K_{g\ell}z^t.
  \label{eq:structural-feedback}
\end{equation}

이 항은 reduced block derivation에서 실제로 nonzero일 때만 사용한다.
Cross blocks가 0인 basis라면 식~\eqref{eq:structural-feedback}를 억지로 만들어내지 않는다.
Module 5는 이전 Local state로 만든
\(Q_{g\leftarrow\ell}^{\mathrm{predictor},t}\)를 이미 소비하므로,
same-frame corrector는 전체 current cross load가 아니라

\begin{equation}
  \Delta Q_{g\leftarrow\ell}^{\mathrm{struct},t}
  =
  Q_{g\leftarrow\ell}^{\mathrm{current},t}
  -Q_{g\leftarrow\ell}^{\mathrm{predictor},t}
\end{equation}

만 소비한다. \(z^t=z^{t-1}\)이고 속도도 같으면 이 delta는 정확히 0이어야 한다.

\subsection*{Corrector solve}

\begin{equation}
  \delta q^t
  =
  \operatorname{SolveGlobal}
  \left(
  \Delta Q_g^{\mathrm{aero},t}
  +\Delta Q_{g\leftarrow\ell}^{\mathrm{struct},t}
  \right),
\end{equation}

\begin{equation}
  q^t=\widetilde q^t+\delta q^t.
\end{equation}

한 번의 corrector로 충분한지 실험하고,
추가 fixed-point iteration은 quality--latency ablation으로 둔다.

\subsection*{출력과 다음 단계}

Corrected \(q^t,\dot q^t,z^t,\dot z^t\)가 Module 12로 넘어간다.
One-way baseline은 aerodynamic delta와 structural coupling을 모두 끈다.

\subsection*{필수 검증}

\begin{itemize}[leftmargin=1.5em]
  \item Feedback을 껐을 때 Global modal amplitude/phase가 실제로 달라지는가?
  \item Corrector가 base force를 두 번 적분하지 않는가?
  \item Same-frame corrector와 next-frame-only feedback의 비용/품질 비교.
  \item Feedback이 미미하면 two-way physics를 main claim에서 내릴 것.
\end{itemize}

\section{Module 12: 최종 anchor state와 deformation field}

\subsection*{푸는 질문}

\emph{Corrected Global/Local coordinate를 최종 anchor 위치, 속도,
deformation gradient로 어떻게 한 번만 복원하는가?}

\subsection*{Solver 분류}

\deterministic Basis lifting과 surface-aware affine reconstruction이다.

\subsection*{최종 상태}

\begin{equation}
  x^t=x^0+\Phi q^t+\Psi_{\mathcal S_t}z^t,
  \qquad
  v^t=\Phi\dot q^t+\Psi_{\mathcal S_t}\dot z^t.
\end{equation}

각 anchor/sample의 final deformation gradient \(F_k^t\)는
rest/current neighborhoods에서 한 번만 계산한다.
Global map과 Local map을 Gaussian에 순차 적용하여 double deformation을 만들지 않는다.

\subsection*{출력과 다음 단계}

\begin{equation}
  \mathcal X^t
  =
  \{x_k^t,v_k^t,F_k^t,n_k^t,A_k^t\}_{k=1}^K.
\end{equation}

이 결과는 Module 13의 Gaussian transport와
다음 frame Module 2의 current surface state에 사용된다.

\subsection*{필수 검증}

\begin{itemize}[leftmargin=1.5em]
  \item Global-only, Local-only, combined reconstruction의 linearity/consistency.
  \item Pin/soft attachment 오차.
  \item \(F\) orientation과 determinant/rank degeneracy.
  \item Global-span leakage와 patch seam diagnostics.
\end{itemize}

\section{Module 13: Gaussian affine transport와 rendering}

\subsection*{푸는 질문}

\emph{Sparse anchor의 최종 deformation을 수많은 render Gaussians에
center-only artifact 없이 어떻게 전달하는가?}

\subsection*{Solver 분류}

\deterministic Affine MLS, covariance transport, co-rotation,
standard 3DGS rendering이다.
이 연산 자체는 기존 Gaussian physics 계열의 선행이 있으므로 contribution으로 주장하지 않는다
\cite{xie2024physgaussian,yan2026gaussianswaying,shao2025gausim}.

\subsection*{Gaussian mean}

Gaussian \(j\)에 대해 local affine map \((F_j^t,t_j^t)\)를
binding support \(\mathcal B_j\)에서 맞춘다.

\begin{equation}
  \mu_j^t
  =
  F_j^t\mu_j^0+t_j^t.
\end{equation}

\subsection*{Covariance}

\begin{equation}
  \Sigma_j^t
  =
  F_j^t\Sigma_j^0F_j^{t\top}.
  \label{eq:covariance-transport}
\end{equation}

Eigenvalue clamp 또는 polar decomposition으로
\(\Sigma_j^t\)가 symmetric positive definite를 유지하도록 한다.
Center-only transport는 anisotropic splat orientation과 scale을 보존하지 못하므로 baseline으로 둔다.

\subsection*{Appearance frame}

Canonical appearance SH coefficient를 유지하고
rigid/co-rotational frame change가 필요하면 analytic SH rotation 또는
local view direction transform을 사용한다.
새 appearance residual network는 본 논문의 main contribution에 포함하지 않는다.

\subsection*{출력}

\begin{equation}
  G^t
  =
  \{(\mu_j^t,\Sigma_j^t,\alpha_j,c_j^t)\}_{j=1}^{N_G},
  \qquad
  I_t=\mathcal R_{\mathrm{3DGS}}(G^t,\Pi_t).
\end{equation}

\subsection*{필수 검증}

\begin{itemize}[leftmargin=1.5em]
  \item Identity, translation, rigid rotation, affine reproduction.
  \item Invalid covariance와 scale explosion count.
  \item Center-only, rigid-only, full affine transport 비교.
  \item Novel-view PSNR/SSIM/LPIPS와 temporal flicker.
\end{itemize}

\section{Module 14: State, hysteresis, decay, diagnostics 갱신}

\subsection*{푸는 질문}

\emph{어떤 Local state를 다음 frame까지 유지하고,
언제 expensive network를 끄며, 언제 완전히 dormant로 만들 것인가?}

\subsection*{Solver 분류}

\deterministic State machine과 energy-based cleanup이다.

\subsection*{세 가지 patch 상태}

\begin{longtable}{L{0.18\linewidth}L{0.30\linewidth}L{0.44\linewidth}}
\toprule
상태 & 실행 & 전환 규칙 \\
\midrule
\endfirsthead
\toprule
상태 & 실행 & 전환 규칙 \\
\midrule
\endhead
Dormant & Gate만 평가, Local network/physics allocation 없음 & Benefit가 \(\tau_{\mathrm{on}}\)을 넘으면 Active \\
\addlinespace
Active & Missing-force network, local excitation, local physical solve & Benefit가 \(\tau_{\mathrm{off}}\) 아래로 충분히 유지되면 Decay \\
\addlinespace
Decay & Network off, excitation fade-out, restoring/damping solve 유지 & Local energy가 \(\tau_E\) 아래면 Dormant; benefit가 다시 커지면 Active \\
\bottomrule
\end{longtable}

\subsection*{Energy criterion}

Patch/group local energy를

\begin{equation}
  E_r^t
  =
  \frac12\dot z_r^{t\top}M_r\dot z_r^t
  +
  \frac12 z_r^{t\top}K_rz_r^t
\end{equation}

로 추적한다.
\(E_r^t<\tau_E\)이고 activation score도 낮을 때만 dormant로 만든다.
Dormant 전환 시 이미 수치적으로 작은 state만 제거하므로 popping을 제한한다.

\subsection*{저장 항목}

\begin{itemize}[leftmargin=1.5em]
  \item \(q^t,\dot q^t,z^t,\dot z^t\),
  \item active/decay/dormant state와 \(\gamma_r^t\),
  \item recurrent aerodynamic memory \(h_r^t\)가 있다면 그 state,
  \item force/torque/energy/conservation residual,
  \item active ratio, network calls, per-module p50/p95 latency.
\end{itemize}

\subsection*{필수 검증}

\begin{itemize}[leftmargin=1.5em]
  \item Threshold-crossing wind에서 chattering과 jerk.
  \item Wind-stop 후 Global/Local energy monotonic decay.
  \item Patch off 순간 position/velocity discontinuity.
  \item Memory leak와 active-state 누적.
\end{itemize}

\section{통합 runtime algorithm}

\begin{lstlisting}
Input:
  static 3DGS G0
  target object package P(G0)
  prescribed spatial wind W(x,t)
  material and attachment controls

State:
  global q, qdot
  local z, zdot over active/decay complementary coordinates
  gate state, activation blend, optional aerodynamic memory

for frame t:
    # Known state and current geometry
    load and validate previous state
    lift q,z to anchor positions and velocities
    fit current surface deformation maps, normals, and areas

    # Global physical predictor
    interpolate wind and compute relative wind
    evaluate current-surface quasi-steady aerodynamic force
    apply conservative cubature and project to Global load
    solve always-on Global reduced structural dynamics -> q_tilde

    # Adaptive Local decision and proposal
    evaluate lightweight future-benefit gate on all patches
    select active patches under compute budget and hysteresis
    evaluate analytic Local base aero on active/halo samples
    evaluate missing-force network only on active patches

    # Assemble before integrating
    keep base aero, missing aero, missing structural as separate channels
    scatter each channel to shared anchors and solve its conservation constraints
    remove/route each channel's Global-span force components
    build one active/decay complementary Local system
    integrate Local mass-stiffness-damping system -> z

    # Mechanical feedback
    reconstruct corrected surface in active/halo regions
    recompute only corrected-minus-predictor aerodynamic force
    add structural_cross_delta = updated coupling - predictor-consumed coupling
    solve one small Global corrector -> q

    # Output and state
    lift final q,z once to anchors
    transport Gaussian means and covariances by affine MLS
    render standard 3DGS
    update active/decay/dormant state and diagnostics
\end{lstlisting}

\subsection{One-frame data dependency}

\begin{equation}
\boxed{
\begin{aligned}
s^{t-1},W^t
&\rightarrow
\text{current surface}
\rightarrow
f_{\mathrm{aero}}^t
\rightarrow
\widetilde q^t\\
&\rightarrow
\text{benefit gate}
\rightarrow
\text{active patch Local base + missing force channels}\\
&\rightarrow
\text{conservative complement assembly}
\rightarrow
z^t\\
&\rightarrow
\Delta Q_g^{\mathrm{aero}},\Delta Q_g^{\mathrm{structural\ cross}}
\rightarrow
q^t
\rightarrow
G^t,s^t.
\end{aligned}
}
\end{equation}

\section{학습 objective}

\subsection{Topology와 physical package loss}

Candidate material edge classification은 class imbalance를 고려한다.

\begin{equation}
  \mathcal L_{\mathrm{edge}}
  =
  \operatorname{BCE}_{\mathrm{focal}}
  (\widehat y_{ij}^{\mathrm{mat}},y_{ij}^{*,\mathrm{mat}}).
\end{equation}

Rest length, tangent direction, oriented normal, area, mass는

\begin{equation}
\begin{aligned}
  \mathcal L_{\mathrm{geom}}
  =&\;
  \lambda_\ell\|\widehat\ell^0-\ell^{0,*}\|_1
  +\lambda_t(1-|\widehat t^\top t^*|)\\
  &+\lambda_n(1-\widehat n^\top n^*)
  +\lambda_A\|\widehat A-A^*\|_1
  +\lambda_m\|\widehat m-m^*\|_1.
\end{aligned}
\end{equation}

Total mass/area conservation은 별도로 검사한다.

\begin{equation}
  \mathcal L_{\mathrm{total}}
  =
  \left|
  \sum_k\widehat m_k-M^*
  \right|
  +
  \left|
  \sum_k\widehat A_k-A^*
  \right|.
\end{equation}

Attachment는 classification과 compliance regression을 분리한다.
Disconnected close layers를 잘못 연결하는 edge에 큰 penalty를 둔다.

\subsection{Reduced operator와 modal loss}

Predicted physical package에서 얻은 eigenfrequency와 mode subspace를 teacher와 비교한다.
Mode의 sign과 repeated eigenvalue ambiguity 때문에 개별 vector L2보다 subspace projector를 사용한다.

\begin{equation}
  \mathcal L_{\mathrm{subspace}}
  =
  \left\|
  \Phi\Phi^\top M
  -
  \Phi^*\Phi^{*\top}M^*
  \right\|_F^2.
\end{equation}

\begin{equation}
  \mathcal L_{\mathrm{freq}}
  =
  \sum_{j=1}^{r_g}
  \left|
  \log(\widehat\omega_j+\eps)
  -
  \log(\omega_j^*+\eps)
  \right|.
\end{equation}

Operator construction은 학습 loss와 별도로 다음 hard checks를 통과해야 한다.

\begin{equation}
  M\succ0,\qquad
  K\succeq0,\qquad
  C\succeq0,\qquad
  K_{g\ell}=K_{\ell g}^\top,\qquad
  C_{g\ell}=C_{\ell g}^\top.
\end{equation}

\subsection{Aerodynamic cubature loss}

식~\eqref{eq:cubature-fit}의 generalized force와 wrench loss에
held-out deformation regularization을 추가한다.

\begin{equation}
  \mathcal L_{\mathrm{cub}}
  =
  \lambda_Q\mathcal L_Q
  +\lambda_F\mathcal L_F
  +\lambda_\tau\mathcal L_\tau
  +\lambda_{\mathrm{sparse}}|\mathcal S_c|.
\end{equation}

단순히 training force snapshot을 맞추고
current normal/area 또는 traveling gust에 실패하면 hyper-reduction claim을 지지하지 못한다.

\subsection{Gate loss}

Continuous benefit regression은

\begin{equation}
  \mathcal L_{\mathrm{benefit}}
  =
  \operatorname{Huber}
  (\widehat b_r^t,b_r^{*,t}).
\end{equation}

같은 frame의 patch 순위를 맞추기 위한 ranking loss는

\begin{equation}
  \mathcal L_{\mathrm{rank}}
  =
  \sum_{(r,s):b_r^*>b_s^*}
  \max
  \left(
    0,\;
    m-\widehat b_r+\widehat b_s
  \right).
\end{equation}

Uncertainty calibration과 budget-aware objective를 더한다.

\begin{equation}
  \mathcal L_{\mathrm{gate}}
  =
  \lambda_b\mathcal L_{\mathrm{benefit}}
  +\lambda_r\mathcal L_{\mathrm{rank}}
  +\lambda_{\mathrm{cal}}\mathcal L_{\mathrm{calibration}}
  +\lambda_c\mathcal L_{\mathrm{cost}}.
\end{equation}

Compute regularizer를 너무 일찍 켜면 모든 patch를 끄는 collapse가 생길 수 있다.
먼저 benefit prediction과 oracle recall을 학습한 뒤 cost term을 활성화한다.

\subsection{Missing-force loss}

Force-space supervision은

\begin{equation}
  \mathcal L_{\Delta f}
  =
  \sum_{t,r}
  w_r^t
  \operatorname{Huber}
  \left(
    \widehat{\Delta f}_r^t,
    \Delta f_r^{*,t}
  \right).
\end{equation}

Global leakage penalty는

\begin{equation}
  \mathcal L_{\mathrm{leak}}
  =
  \sum_t
  \left\|
  \Phi^\top f_{\ell,\theta}^t
  \right\|_2^2.
\end{equation}

Internal structural channel의 net force/torque loss는

\begin{equation}
  \mathcal L_{\mathrm{int\mbox{-}wrench}}
  =
  \left\|
  \sum_k\Delta f_{k,\mathrm{struct}}
  \right\|_2^2
  +
  \left\|
  \sum_k(x_k-c)\times\Delta f_{k,\mathrm{struct}}
  \right\|_2^2.
\end{equation}

External aerodynamic channel은 0으로 만들지 않고 teacher external wrench를 맞춘다.

\subsection{One-step과 rollout state loss}

\begin{equation}
  \mathcal L_{\mathrm{state}}
  =
  \lambda_q\|q-q^*\|_1
  +\lambda_{\dot q}\|\dot q-\dot q^*\|_1
  +\lambda_z\|z-z^*\|_1
  +\lambda_{\dot z}\|\dot z-\dot z^*\|_1.
\end{equation}

Multi-step rollout은

\begin{equation}
  \mathcal L_{\mathrm{roll}}
  =
  \sum_{h=1}^{H}
  \beta^{h-1}
  \left[
  \lambda_x\|x^{t+h}-x^{*,t+h}\|_1
  +\lambda_v\|v^{t+h}-v^{*,t+h}\|_1
  +\lambda_n\mathcal L_n^{t+h}
  \right].
\end{equation}

Flutter의 amplitude만 맞고 frequency/phase가 틀리는 것을 막기 위해
window spectrum loss를 사용한다.

\begin{equation}
  \mathcal L_{\mathrm{spec}}
  =
  \left\|
  \log(\operatorname{PSD}(\widehat x)+\eps)
  -
  \log(\operatorname{PSD}(x^*)+\eps)
  \right\|_1.
\end{equation}

\subsection{Overlap, feedback, rendering loss}

\begin{equation}
  \mathcal L_{\mathrm{overlap}}
  =
  \lambda_c\|\Delta x_{\mathrm{boundary}}\|_2^2
  +\lambda_v\|\Delta v_{\mathrm{boundary}}\|_2^2
  +\lambda_w\|A_{\mathrm{cons}}f-b_{\mathrm{cons}}\|_2^2.
\end{equation}

Corrected-aero feedback은 teacher Global generalized force delta와 비교한다.

\begin{equation}
  \mathcal L_{\mathrm{feedback}}
  =
  \left\|
  \Delta Q_g^{\mathrm{aero}}
  -
  \Delta Q_g^{\mathrm{aero},*}
  \right\|_1.
\end{equation}

Gaussian transport는 mean/covariance와 rendered frame에서 감독할 수 있다.

\begin{equation}
  \mathcal L_{\Sigma}
  =
  \left\|
  \log\Sigma_j^t-\log\Sigma_j^{*,t}
  \right\|_F^2,
\end{equation}

\begin{equation}
  \mathcal L_{\mathrm{render}}
  =
  \lambda_1\|\widehat I-I^*\|_1
  +\lambda_s(1-\operatorname{SSIM}(\widehat I,I^*))
  +\lambda_p\operatorname{LPIPS}(\widehat I,I^*).
\end{equation}

\subsection{Total objective와 staged activation}

\begin{equation}
\begin{aligned}
  \mathcal L
  =&\;
  \mathcal L_{\mathrm{topology/geometry}}
  +\mathcal L_{\mathrm{operator/modal}}
  +\mathcal L_{\mathrm{cub}}\\
  &+\mathcal L_{\mathrm{gate}}
  +\mathcal L_{\Delta f}
  +\mathcal L_{\mathrm{leak}}
  +\mathcal L_{\mathrm{wrench}}\\
  &+\mathcal L_{\mathrm{state}}
  +\mathcal L_{\mathrm{roll}}
  +\mathcal L_{\mathrm{spec}}
  +\mathcal L_{\mathrm{overlap}}
  +\mathcal L_{\mathrm{feedback}}
  +\mathcal L_{\mathrm{render}}.
\end{aligned}
\end{equation}

모든 항을 처음부터 동시에 켜지 않는다.
Stage별로 학습하고 각 stage의 oracle upper bound를 먼저 확인한다.

\section{학습 curriculum}

\subsection{Stage 0: Deterministic physics/transport unit tests}

Network training 전에 다음을 통과한다.

\begin{itemize}[leftmargin=1.5em]
  \item teacher coordinate와 independent GS coordinate alignment,
  \item conservative area/mass partition,
  \item attachment force/projection ownership,
  \item \(M,K,C\) positivity와 symmetry,
  \item Global/Local \(M\)-orthogonality,
  \item area-normal cofactor update,
  \item affine MLS identity/rigid/affine reproduction,
  \item covariance SPD.
\end{itemize}

\subsection{Stage 1: Oracle topology Global physics}

Teacher mesh topology, mass, stiffness, attachment를 그대로 사용하고
Global reduced physics만 구현한다.
이 단계에서 current-state analytic aero와 full-sample projection을 먼저 사용한다.
Global-only가 teacher의 큰 motion과 wind-stop recovery를 학습 없이 재현해야 한다.

\subsection{Stage 2: Oracle topology Local physics}

Oracle Local basis와 oracle active set을 사용하여
physics-only Local과 full-resolution Local upper bound를 비교한다.
Global rank를 늘리는 것만으로 local flutter가 해결되는지도 이 단계에서 반증한다.

\subsection{Stage 3: Always-on missing-force network}

Gate 없이 모든 candidate patch의 missing-force network를 실행한다.
Network capacity와 label correctness를 먼저 검증한다.
Always-on model이 physics-only Local을 개선하지 못하면 gate 학습으로 넘어가지 않는다.

\subsection{Stage 4: Gate와 adaptive execution}

Oracle-benefit labels로 gate를 학습한다.
초기에는 quality target만 사용하고,
이후 measured GPU cost와 active-budget regularizer를 넣는다.
Hysteresis와 off-policy noisy rollout state를 학습/검증한다.

\subsection{Stage 5: Predicted topology와 operator}

Oracle topology를 learned topology/distilled operator로 교체한다.
Topology error와 dynamics error를 동시에 처음부터 디버깅하지 않는다.
Oracle 대비 degradation을 분리해 보고한다.

\subsection{Stage 6: Independent GS reconstruction robustness}

다른 view 수, random seed, density, holes, floaters를 가진 independent GS variants로
object package 추론과 dynamics를 평가한다.
Source mesh에 Gaussian을 직접 bind한 clean variant만으로 결과를 주장하지 않는다.

\subsection{Stage 7: End-to-end Gaussian rendering}

Dynamics가 안정된 뒤 affine Gaussian transport와 rendering loss를 활성화한다.
Appearance correction은 별도 supporting experiment로 제한한다.

\section{계산 복잡도와 runtime contract}

\subsection{One-time target setup}

\begin{equation}
\begin{aligned}
  T_{\mathrm{setup}}
  =&\;
  T_{\mathrm{GS\mbox{-}preprocess}}
  +T_{\mathrm{anchor}}
  +T_{\mathrm{topology}}\\
  &+T_{\mathrm{operator}}
  +T_{\mathrm{basis}}
  +T_{\mathrm{cubature}}
  +T_{\mathrm{binding}}.
\end{aligned}
\end{equation}

Setup은 target mesh reconstruction을 호출하지 않아야 한다.
Target에서 SuGaR/GOF mesh를 만드는 variant는 end-to-end baseline이지 mainline이 아니다
\cite{guedon2024sugar,yu2024gof}.

\subsection{Per-frame dynamics}

Global mode 수 \(r_g\), active Local coordinate 수 \(r_\ell^A\),
patch 수 \(P\), active patch 수 \(P_A\), cubature sample 수 \(S_c\)라 하면

\begin{equation}
\begin{aligned}
  T_{\mathrm{dyn}}
  \approx&
  T_{\mathrm{surface}}(S_c)
  +T_{\mathrm{aero}}(S_c)
  +T_{\mathrm{global}}(r_g)\\
  &+T_{\mathrm{gate}}(P)
  +T_{\mathrm{net}}(P_A)
  +T_{\mathrm{assembly}}(P_A)
  +T_{\mathrm{local}}(r_\ell^A)
  +T_{\mathrm{corrector}}(r_g).
\end{aligned}
\end{equation}

Linear reduced matrices를 factorize/cache할 수 있는 구간에서는
Global/Local solve가 작은 triangular solve가 된다.
Material/attachment가 frame마다 바뀌면 재조립 비용을 별도로 보고한다.

\subsection{Gaussian transport}

\begin{equation}
  T_{\mathrm{GS\mbox{-}transport}}
  =
  O(N_GK_{\mathrm{MLS}}).
\end{equation}

Dynamics resolution \(K,P,r_g,r_\ell^A\)는 render Gaussian 수 \(N_G\)와 분리되어야 한다.

\subsection{Measured break-even}

Adaptive method의 이론적 sparsity가 실제 GPU speedup을 보장하지 않는다.
작은 patch가 많으면 kernel launch와 gather/scatter가 병목이 될 수 있다.
따라서 target GPU에서

\begin{equation}
  T_{\mathrm{always}},\quad
  T_{\mathrm{adaptive}}(\rho_A),\quad
  T_{\mathrm{physics\mbox{-}only}},\quad
  T_{\mathrm{direct}}
\end{equation}

를 실제 측정한다.
Empirical break-even \(\rho_A^{\mathrm{BE}}\) 이상에서는
always-on batched Local 또는 direct baseline이 더 빠를 수 있다.
이 경우 runtime dispatcher를 두거나 adaptive claim을 축소한다.

\subsection{Runtime mesh-free contract}

\begin{longtable}{L{0.28\linewidth}L{0.25\linewidth}L{0.39\linewidth}}
\toprule
항목 & 허용 여부 & 설명 \\
\midrule
\endfirsthead
\toprule
항목 & 허용 여부 & 설명 \\
\midrule
\endhead
Training source mesh & 허용 & Privileged topology/physics teacher \\
\addlinespace
Target mesh reconstruction & Mainline 금지 & 별도 end-to-end baseline에서만 허용 \\
\addlinespace
Sparse oriented anchor graph & 허용 & Learned material scaffold; watertight triangle mesh 요구 없음 \\
\addlinespace
Runtime MPM/FEM grid & 금지 & Teacher와 full-physics baseline에서만 사용 \\
\addlinespace
Prescribed wind grid/query & 허용 & Fluid dynamics state가 아니라 control input \\
\addlinespace
Runtime CFD/LBM & 금지 & Teacher-FSI 생성 또는 accuracy upper bound \\
\addlinespace
Static GS preprocessing & 허용 & Reliability, anchors, topology, binding inference \\
\bottomrule
\end{longtable}

\section{관련 연구: 유사점, 차이점, 비교 용도}

\subsection{직접적인 target-domain 경계}

\begin{longtable}{L{0.16\linewidth}L{0.23\linewidth}L{0.30\linewidth}L{0.23\linewidth}}
\toprule
방법 & 유사점 & 핵심 차이 & 실험에서의 역할 \\
\midrule
\endfirsthead
\toprule
방법 & 유사점 & 핵심 차이 & 실험에서의 역할 \\
\midrule
\endhead
Gaussian Swaying \cite{yan2026gaussianswaying}
& Current Gaussian surface normal, area, relative wind로 drag/lift 계열 force 계산
& Full MPM dynamics, 수동 simulation/material setup; object ROM과 error-gated Local 없음
& 동일 prescribed wind에서 Gaussian-surface aero front-end를 맞추고 MPM 대 ROM 정확도/latency 비교 \\
\addlinespace
DiffWind \cite{lei2026diffwind}
& Spatial wind, 3DGS-derived particles, wind-driven forward dynamics
& Grid/MPM/LBM 계열; surface-force-to-ROM과 sparse Local complement가 아님
& Spatial-wind subset에서 조건부 end-to-end 비교; teacher/upper-bound 후보 \\
\addlinespace
PhysGaussian \cite{xie2024physgaussian}
& Gaussian/particle deformation gradient와 covariance transport
& MPM 기반 generic material dynamics; wind-specific topology/gate 없음
& 동일 external load subset의 physics cost와 GS transport sanity comparison \\
\addlinespace
DynamicTree \cite{li2026dynamictree}
& Static real 3DGS, compact spectral/modal-like state, external-force response
& Tree/mesh/sparse voxel spectrum; current-surface aero와 Local corrector 없음
& Tree/leaf subset의 Global response와 latency 조건부 비교 \\
\addlinespace
FreeForm \cite{xiang2026freeform}
& Point/GS representation에서 mesh-free reduced deformable simulation
& Wind/thin-surface topology/current aero/error-gated local 없음
& 동일 anchor/외력에서 Global ROM 정확도와 basis construction 비교 \\
\addlinespace
PhysSkin \cite{lei2026physskin}
& Mesh-free/discretization-agnostic reduced skinning field, GS 적용
& Wind-specific force/operator와 local adaptive correction 없음
& Reduced lifting/geometry generalization subset 비교; 주로 구조 비교 \\
\addlinespace
GausSim \cite{shao2025gausim}
& Gaussian kernels의 hierarchical physics-aware dynamics와 계산 절감
& Fixed hierarchy/coarse-to-fine; modal complement와 posterior-error active set이 아님
& 공통 force/input을 만들 수 있을 때만 conditional quantitative comparison \\
\addlinespace
NGFF \cite{li2026ngff}
& Gaussian object graph의 global/local neural force와 ODE dynamics
& Global은 rigid 6-DoF 중심이며 local stress; elastic modal complement와 wind-specific gate가 아님
& Global/local Gaussian force-field capability 비교; task 불일치 시 정성/인용 \\
\addlinespace
PGRD \cite{patel2026pgrd}
& Physics backbone + learned residual, paper/flag, closed-loop rollout, GS rendering
& Full particle graph의 all-particle residual; spatial error gate와 modal complement 없음
& PGRD-style always-on residual 대 adaptive Local의 핵심 정량 비교 \\
\bottomrule
\end{longtable}

\subsection{Global aerodynamic/reduced dynamics 경계}

\begin{longtable}{L{0.16\linewidth}L{0.23\linewidth}L{0.30\linewidth}L{0.23\linewidth}}
\toprule
방법 & 유사점 & 핵심 차이 & 실험에서의 역할 \\
\midrule
\endfirsthead
\toprule
방법 & 유사점 & 핵심 차이 & 실험에서의 역할 \\
\midrule
\endhead
Wind Projection Basis \cite{diener2009windbasis}
& Wind load를 object modes에 투영하고 reduced dynamics 적분
& Rest geometry, uniform/directional wind 가정이 중심; current-deformed GS surface 없음
& Matched WPB-style baseline으로 current normal/area, relative/spatial wind의 필요성 검증 \\
\addlinespace
OmniAD \cite{martin2015omniad}
& Directional aerodynamic force/torque를 spherical harmonics로 표현
& Rigid object 전체, deformation state 없음
& RSH novelty 경계 인용; analytic/LUT/MLP/RSH 공력 표현 비교 \\
\addlinespace
MeshGraph\-Nets \cite{pfaff2021meshgraphnets}
& Current mesh state와 wind에서 flag-like dynamics를 direct neural rollout
& Mesh 입력/all-node graph network; explicit reduced physics base 없음
& 동일 topology/data로 direct neural Global 일반화와 안정성 비교 \\
\addlinespace
Nonlinear subspace integration \cite{barbic2005stvk}
& Reduced second-order nonlinear structural dynamics
& Generic volumetric elasticity; GS/wind/local gate 없음
& 수식의 선행 경계와 Global solver 구현 참고 \\
\addlinespace
Optimized cubature \cite{an2008cubature}
& Subspace force integration을 sparse samples로 가속
& GS surface, spatial aero, net wrench 목적이 아님
& Cubature 선행 인용; 본 연구는 force/moment-preserving aero adaptation을 검증 \\
\bottomrule
\end{longtable}

\subsection{Local correction/adaptivity 경계}

\begin{longtable}{L{0.16\linewidth}L{0.23\linewidth}L{0.30\linewidth}L{0.23\linewidth}}
\toprule
방법 & 유사점 & 핵심 차이 & 실험에서의 역할 \\
\midrule
\endfirsthead
\toprule
방법 & 유사점 & 핵심 차이 & 실험에서의 역할 \\
\midrule
\endhead
Subspace Condensation \cite{teng2015condensation}
& Novel event 주변 full-space region을 runtime 활성화, Global/Local coupling
& Volumetric contact, analytic adaptive solve; learned benefit/missing force/GS wind 없음
& Adaptive-local 선행 경계 인용; 동일 구현 비교보다 conceptual comparison \\
\addlinespace
Trading Spaces \cite{trusty2024tradingspaces}
& Solution quality oracle, local enrichment, quality--cost adaptivity
& Full-space contact elastodynamics; static GS/wind/neural defect 없음
& Gate/oracle 설계의 선행 경계; oracle benefit 및 tolerance curve 비교 개념 \\
\addlinespace
Learning Contact Corrections \cite{romero2021contact}
& Global reduced deformation 위 learned high-frequency local correction
& Contact-specific, always evaluated correction; posterior spatial gate와 wind feedback 없음
& Direct local displacement correction ablation과 선행 인용 \\
\addlinespace
Fast Complementary Dynamics \cite{benchekroun2023complementary}
& Primary motion과 겹치지 않는 complementary secondary dynamics
& Rig-driven secondary motion; error-gated patch missing force와 wind 없음
& \(M\)-orthogonal complement 선행 경계와 no-complement ablation의 근거 \\
\addlinespace
Physics-Inspired Cloth Upsampling \cite{kavan2011upsampling}
& Coarse cloth 위 high-frequency learned/oscillatory detail, wind flag 사례
& Offline/detail upsampling 중심; current aero와 closed-loop local force 아님
& Global motion+detail 개념 경계 인용; 직접 quantitative baseline은 선택 \\
\bottomrule
\end{longtable}

\subsection{비교 구현의 명명 규칙}

공식 코드와 원 설정을 재현하지 못한 경우 논문명을 그대로 baseline 이름으로 쓰지 않는다.
다음처럼 구현 범위를 명시한다.

\begin{itemize}[leftmargin=1.5em]
  \item \emph{Wind Projection Basis-style rest-state modal load}.
  \item \emph{Gaussian Swaying-style surface aero + matched MPM}.
  \item \emph{PGRD-style all-node physics residual}.
  \item \emph{MeshGraphNets-style direct graph rollout}.
\end{itemize}

입력, topology, DOF, timestep, wind control을 맞추지 못한 방법은
main quantitative leaderboard가 아니라 capability/failure-mode comparison으로 분리한다.

\clearpage
\section{비교 실험 대장}

이 절은 방법 설명에서 분리된 실험 계획이다.
모든 실험은 먼저 검증하려는 claim을 정하고,
통제 변수, 비교군, 지표, 실패 시 결정을 명시한다.
모든 조합의 Cartesian product를 실행하기보다
단계별 controlled test와 최종 모델의 one-factor ablation을 병행한다.

\subsection{검증할 핵심 claim}

\begin{longtable}{L{0.07\linewidth}L{0.37\linewidth}L{0.48\linewidth}}
\toprule
ID & Claim & 반드시 필요한 증거 \\
\midrule
\endfirsthead
\toprule
ID & Claim & 반드시 필요한 증거 \\
\midrule
\endhead
C1 & Training-only mesh에서 distill한 topology/operator가 target runtime mesh 없이 thin-surface dynamics를 유지한다.
& Oracle topology, predicted topology, kNN/FPS, target mesh reconstruction baseline의 topology와 downstream dynamics 차이 \\
\addlinespace
C2 & Global reduced physics가 object-wide 저주파 response를 안정적이고 controllable하게 계산한다.
& Full physics, sparse physics, WPB-style, direct/recurrent neural Global, hybrid Global과의 trajectory/frequency/damping/latency 비교 \\
\addlinespace
C3 & Current-state spatial aerodynamic hyper-reduction이 적은 sample로 generalized force와 net wrench를 보존한다.
& Full-sample, rest-state, random/FPS, nonconservative learned samples, proposed conservative samples 비교 \\
\addlinespace
C4 & Local physics + learned missing force가 Global이 놓친 flutter/fold를 복원한다.
& Global-only, full Local physics, reduced Local physics, hybrid missing-force, direct displacement residual 비교 \\
\addlinespace
C5 & Learned future-benefit gate가 always-on Local과 유사한 품질을 더 낮은 실제 비용으로 달성한다.
& Never, Always, Random, Heuristic, Oracle, Learned gate의 동일 budget 및 동일 latency 비교 \\
\addlinespace
C6 & Complement, overlap conservation, decay, feedback이 double counting과 비물리적 seam/drift를 방지한다.
& 각 coupling 요소 제거 ablation과 force/torque/energy/jerk/Global leakage 측정 \\
\addlinespace
C7 & 전체 시스템이 wind-driven static 3DGS에서 유용한 quality--latency Pareto와 OOD/장기 안정성을 제공한다.
& Matched solver track, end-to-end static-GS track, real captured GS demo, 30--60초 stress test \\
\bottomrule
\end{longtable}

\subsection{평가 Track 분리}

\subsubsection{Track A: Oracle-topology matched solver comparison}

모든 방법에 다음 정보를 동일하게 제공한다.

\begin{itemize}[leftmargin=1.5em]
  \item 동일 teacher surface/topology,
  \item 동일 mass, material, attachment,
  \item 동일 wind field와 initial state,
  \item 동일 timestep,
  \item 가능한 한 동일 anchor 수/DOF,
  \item 동일 rendering transport.
\end{itemize}

이 Track은 topology reconstruction 품질을 배제하고
aero operator, Global solver, Local correction, gate, coupling, runtime만 비교한다.
Method accuracy의 주장은 Track A에서 먼저 성립해야 한다.

\subsubsection{Track B: Static-3DGS end-to-end comparison}

각 방법은 원래 요구하는 입력을 사용하되, 입력 특권을 명시한다.

\begin{itemize}[leftmargin=1.5em]
  \item 제안 방법: independently reconstructed static 3DGS, material/attachment controls.
  \item Target mesh baseline: target static GS에서 mesh를 별도 추출한 뒤 simulation.
  \item MPM baseline: 해당 preprocessing과 particle/grid representation.
  \item Neural baseline: 요구하는 topology/mesh/data를 표에 명시.
\end{itemize}

이 Track은 target-mesh-free property, preprocessing robustness,
최종 dynamics/rendering, setup/runtime cost를 평가한다.
입력 특권이 다른 결과를 하나의 공정한 solver leaderboard로 해석하지 않는다.

\subsubsection{Track C: Real captured static GS}

실제 multi-view capture에서 재구성한 flag/cloth/leaf asset을 사용한다.
가능하면 marker, silhouette, tracked feature, wind sensor를 수집한다.
정확한 material/force ground truth가 없으면 정성 결과와 제한된 trajectory/silhouette metric으로 보고한다.

\subsection{Test asset와 wind program}

\begin{longtable}{L{0.19\linewidth}L{0.25\linewidth}L{0.25\linewidth}L{0.23\linewidth}}
\toprule
Asset & 핵심 현상 & Attachment & 주요 실험 \\
\midrule
\endfirsthead
\toprule
Asset & 핵심 현상 & Attachment & 주요 실험 \\
\midrule
\endhead
Short cantilever strip & 단일 bending mode, damping, 시작/정지 & 한쪽 edge hard pin & Global unit test, frequency/damping \\
\addlinespace
Long flag & traveling wave, tip flutter, mode truncation & 한쪽 edge hard/compliant & Global/Local/gate main benchmark \\
\addlinespace
Pennant & 비대칭 geometry, tip motion & 한쪽 short edge & Generalization과 aero torque \\
\addlinespace
Hanging cloth & 여러 attachment와 broad deformation & top edge 또는 two points & Topology/operator, Global capacity \\
\addlinespace
Curtain strip & spatially varying gust와 long-range response & top edge/track & Spatial wind, active patch 이동 \\
\addlinespace
Paper sheet & stiff bending, fold-like local response & two points/partial clamp & Local physics/missing-force \\
\addlinespace
Single leaf & anisotropy, stem attachment, twist & compliant stem & Material / attachment OOD \\
\addlinespace
Simple multi-leaf plant & 여러 component와 batch behavior & stems & End-to-end GS scaling; full tree claim은 하지 않음 \\
\bottomrule
\end{longtable}

Wind program은 다음을 포함한다.

\begin{enumerate}[leftmargin=1.7em]
  \item Zero-wind relaxation.
  \item Constant weak/medium/strong wind.
  \item Abrupt wind start와 abrupt stop.
  \item Sinusoidal gust.
  \item Frequency chirp.
  \item Yaw/elevation directional sweep.
  \item Traveling gust.
  \item Localized spatial gust.
  \item OOD strong wind failure test.
\end{enumerate}

\subsection{공통 baseline registry}

\begin{longtable}{L{0.08\linewidth}L{0.30\linewidth}L{0.29\linewidth}L{0.25\linewidth}}
\toprule
ID & Baseline & 통제/입력 & 검증 목적 \\
\midrule
\endfirsthead
\toprule
ID & Baseline & 통제/입력 & 검증 목적 \\
\midrule
\endhead
B0 & High-resolution nonlinear thin-shell/MPM teacher & Full mesh/full resolution & Accuracy upper bound \\
\addlinespace
B1 & Same-anchor sparse corotational shell, PD, 또는 XPBD & 동일 graph/DOF/wind & MPM 대비 해상도 차이만으로 얻은 speedup인지 확인 \\
\addlinespace
B2 & WPB-style rest-state modal wind load & 동일 \(\Phi,M,C,K\), rest normal/area, uniform wind & Current state/spatial aero contribution \\
\addlinespace
B3 & Current aero + all samples + Global ROM & Oracle topology, no hyper-reduction & Global aero accuracy upper bound \\
\addlinespace
B4 & Direct neural Global MLP/GNN & 동일 state/features/data & Physics Global 필요성 \\
\addlinespace
B5 & Recurrent neural Global GRU/state-space & 동일 state/history/data & Neural memory를 포함한 공정 비교 \\
\addlinespace
B6 & Physics Global + learned \(\Delta Q_g\) & 동일 Global base & Global hybrid correction 필요성 \\
\addlinespace
B7 & Global-only proposed & Local off & Local 순수 기여 \\
\addlinespace
B8 & Global + full-resolution Local physics & Oracle active region/full local solver & Local accuracy upper bound \\
\addlinespace
B9 & Global + reduced Local physics only & 동일 Local basis, network off & Missing-force network 필요성 \\
\addlinespace
B10 & Global + direct local displacement network & 동일 patch features/data & Force prediction+integration 필요성 \\
\addlinespace
B11 & PGRD-style all-node/all-patch residual & 동일 base graph와 training data & Always-on physics residual 대 adaptive Local \\
\addlinespace
B12 & Proposed + always-on Local expert & Gate off, all patches active & Learned Local quality upper bound/비용 상한 \\
\addlinespace
B13 & Proposed + oracle-benefit gate & Offline oracle selection & Gate selection upper bound \\
\addlinespace
B14 & Proposed full learned gate & 최종 method & Main result \\
\bottomrule
\end{longtable}

\section{E0: Deterministic contract와 단위 테스트}

\subsection*{목적}

Network error와 physics/geometry implementation bug를 분리한다.
한 항목이라도 실패하면 long training을 시작하지 않는다.

\subsection*{테스트}

\begin{longtable}{L{0.26\linewidth}L{0.34\linewidth}L{0.32\linewidth}}
\toprule
대상 & 입력 & 통과 조건 \\
\midrule
\endfirsthead
\toprule
대상 & 입력 & 통과 조건 \\
\midrule
\endhead
Coordinate alignment & Known mesh/GS cameras와 marker & Metric scale/axis/handedness 일치 \\
\addlinespace
Area/mass partition & Dense Gaussian support & Total area/mass 보존, density dependence 없음 \\
\addlinespace
Operator & Random valid materials & \(M\succ0,K,C\succeq0\), symmetry, attachment null-space 처리 \\
\addlinespace
Modal complement & Random \(q,z\) & \(\|\Phi^\top M\Psi\|\) tolerance 이하 \\
\addlinespace
Aero & Zero/rigid/relative wind & Zero relative wind force 0, rotation equivariance \\
\addlinespace
Cubature & Held-out force snapshots & Generalized force와 wrench tolerance \\
\addlinespace
Overlap & Synthetic opposing patch forces & Constraint residual, net force/torque 보존 \\
\addlinespace
Integrator & Zero wind and known oscillator & Frequency/damping/energy behavior 일치 \\
\addlinespace
Corrector & Synthetic Local deformation & Corrected-minus-predictor delta만 적용 \\
\addlinespace
GS transport & Identity, translation, rotation, affine & Mean / covariance reproduction, SPD 유지 \\
\bottomrule
\end{longtable}

\section{E1: Topology와 operator distillation}

\subsection*{검증 claim}

C1: Static GS에서 만든 target-runtime package가 Oracle topology에 가까운 dynamics를 제공하는가?

\subsection*{비교군}

\begin{enumerate}[leftmargin=1.7em]
  \item Oracle source mesh topology/operator.
  \item Proposed oriented topology distillation.
  \item Euclidean kNN graph.
  \item FPS anchors + kNN.
  \item Orientation/bending relation 제거.
  \item Target GS에서 SuGaR/GOF-style mesh extraction 후 physics.
  \item Attachment oracle 대 predicted/user-provided attachment.
\end{enumerate}

\subsection*{정적 topology 지표}

\begin{itemize}[leftmargin=1.5em]
  \item Material edge precision/recall/F1.
  \item Close disconnected layer false-positive rate.
  \item Normal orientation consistency.
  \item Attachment classification/calibration.
  \item Area/mass relative error.
  \item Eigenfrequency relative error.
  \item Modal subspace principal angles.
\end{itemize}

\subsection*{Downstream 지표}

Global/Local trajectory, normal/curvature, aerodynamic load,
long-rollout stability를 Oracle topology와 비교한다.
Topology F1만 높고 downstream dynamics가 나쁘면 C1은 성립하지 않는다.

\subsection*{통제 변수}

동일 teacher state, wind, material, timestep, Global rank,
Local basis budget, aerodynamic samples를 사용한다.

\subsection*{반증과 결정}

\killgate Predicted topology dynamics error가 Oracle 대비 지속적으로 크게 악화되거나
close-layer false connection이 안정성을 깨면 target-mesh-free claim을 보류한다.
첫 paper를 oracle/simple topology asset으로 축소하거나
target mesh extraction을 명시적인 setup baseline으로 채택한다.

\section{E2: Aerodynamic operator와 RSH 대안}

\subsection*{검증 claim}

C3: Current-deformed surface와 conservative hyper-reduction이 필요한가?
RSH 없이 direct analytic evaluation이 더 나은가?

\subsection*{실험 격리}

Teacher/current surface positions, velocities, normals, areas를 고정 입력으로 제공하고
dynamics integration을 끈다.
이렇게 해야 구조 solver error를 aerodynamic operator error로 오해하지 않는다.

\subsection*{비교군}

\begin{enumerate}[leftmargin=1.7em]
  \item Rest surface + uniform wind modal load: WPB-style.
  \item Current surface + absolute wind.
  \item Current surface + relative wind + all samples.
  \item Current surface + random samples.
  \item Current surface + FPS samples.
  \item Learned samples without wrench constraint.
  \item Proposed conservative hyper-reduction.
  \item Directional LUT + interpolation.
  \item Small directional MLP.
  \item Fitted RSH upper bound.
  \item Predicted/state-conditioned RSH.
  \item Teacher-FSI force, 해당 subset이 있을 때만.
\end{enumerate}

\subsection*{입력 sweep}

Held-out yaw/elevation, wind speed, deformed state,
surface velocity, traveling gust, local gust,
different GS reconstruction을 사용한다.

\subsection*{지표}

\begin{itemize}[leftmargin=1.5em]
  \item Sample force error.
  \item Global generalized-force relative error.
  \item Net force와 net torque relative error.
  \item Drag/lift coefficient error.
  \item Unsteady subset의 amplitude/phase error.
  \item Sample count, memory, p50/p95 GPU latency.
\end{itemize}

\subsection*{결정 규칙}

\killgate Direct current-surface analytic force가 RSH/LUT/MLP보다
정확하고 같거나 더 빠르면 RSH는 method에서 제거하고 comparison 결과로만 남긴다.
Conservative sampling이 random/FPS보다 dynamics/wrench drift에서 이득이 없으면
hyper-reduction을 main contribution으로 주장하지 않는다.

\section{E3: Global physics 대 neural Global}

\subsection*{검증 claim}

C2: Global reduced physics가 direct neural state transition보다
안정성, control, generalization에서 이점이 있는가?

\subsection*{비교군}

\begin{enumerate}[leftmargin=1.7em]
  \item Full teacher.
  \item Same-anchor sparse PD/XPBD/corotational shell.
  \item WPB-style Global.
  \item Proposed physics Global.
  \item Physics Global + learned generalized-force correction.
  \item Direct next-state MLP.
  \item MeshGraphNets-style direct GNN rollout.
  \item Recurrent GRU/state-space Global.
\end{enumerate}

Direct neural baselines는 가능한 한 동일 topology,
동일 wind/material/attachment features, 동일 data volume을 사용한다.

\subsection*{조건}

Local을 모두 끄고 Oracle topology를 사용한다.
Steady wind, start/stop, sinusoidal gust, chirp, traveling gust,
unseen direction/speed/material/attachment에서 평가한다.

\subsection*{지표}

\begin{itemize}[leftmargin=1.5em]
  \item Object-scale normalized position/velocity error.
  \item Tip amplitude와 phase lag.
  \item Dominant frequency와 mode response.
  \item Damping ratio와 wind-stop envelope.
  \item Work--energy residual과 divergence rate.
  \item 4--10초 및 30--60초 rollout.
  \item Per-step p50/p95 latency와 memory.
\end{itemize}

\subsection*{결정 규칙}

Direct neural Global이 OOD와 long rollout에서도 동등하거나 우수하고 더 빠르면
Global physics 선택을 재검토한다.
Physics + learned \(\Delta Q_g\)가 pure physics를 유의미하게 개선하지 않으면
Global neural correction을 제거한다.

\section{E4: Local physics와 missing-force network}

\subsection*{검증 claim}

C4: Global basis가 놓친 detail이 실제로 존재하며,
reduced Local physics와 learned missing force가 각각 필요한가?

\subsection*{비교군}

\begin{longtable}{L{0.30\linewidth}L{0.58\linewidth}}
\toprule
Variant & 검증 내용 \\
\midrule
\endfirsthead
\toprule
Variant & 검증 내용 \\
\midrule
\endhead
Global-only & Local contribution 자체의 필요성 \\
\addlinespace
Global rank 증가 & Local 대신 Global capacity 확대로 해결 가능한지 \\
\addlinespace
Full-resolution Local physics & 비용을 무시한 local accuracy upper bound \\
\addlinespace
Reduced Local physics-only & 저렴한 known physics만으로 충분한지 \\
\addlinespace
Reduced Local + learned missing force & 최종 Local expert의 추가 이득 \\
\addlinespace
Learned force, local integrator 없음 & 물리 적분기의 필요성 \\
\addlinespace
Direct local displacement/velocity network & Force-space hybrid의 안정성 이점 \\
\addlinespace
PGRD-style all-patch residual & Sparse Local 선택의 이점 \\
\bottomrule
\end{longtable}

\subsection*{평가 영역}

Flag free tip, fold/crease 주변, attachment transition,
traveling gust front, topology uncertainty가 큰 영역을 따로 보고한다.

\subsection*{지표}

\begin{itemize}[leftmargin=1.5em]
  \item Patch-local position/velocity/normal/curvature error.
  \item High-frequency power spectrum distance.
  \item Flutter amplitude, frequency, phase.
  \item Wind-stop decay와 residual energy.
  \item Local force error와 local compute time.
  \item Global-span leakage.
\end{itemize}

\subsection*{결정 규칙}

\killgate Global mode 수를 합리적으로 늘리면 Local과 차이가 사라질 경우
Local contribution을 축소한다.
Reduced Local physics-only가 learned missing force와 동등하면 network를 제거하고
pure adaptive physics로 전환한다.
Direct displacement가 같은 데이터에서 더 안정적이고 빠르면
force-space 설계를 다시 검토하되 physical control/OOD 결과를 함께 판단한다.

\section{E5: Gate와 실제 adaptive compute}

\subsection*{검증 claim}

C5: Learned gate가 Local expert의 품질을 유지하면서 실제 실행 비용을 줄이는가?

\subsection*{동일 Local expert를 사용한 gate 비교}

\begin{enumerate}[leftmargin=1.7em]
  \item Never Local.
  \item Always Local.
  \item Random gate with matched active ratio.
  \item Curvature threshold.
  \item Strain/deformation severity threshold.
  \item Teacher residual threshold: runtime 불가능한 분석 기준.
  \item Oracle future-benefit gate.
  \item Learned binary error gate.
  \item Proposed continuous cost-normalized benefit gate.
\end{enumerate}

\subsection*{Gate 자체 지표}

\begin{itemize}[leftmargin=1.5em]
  \item High-benefit patch recall/precision/AUPRC.
  \item Expected calibration error.
  \item Oracle 대비 benefit regret.
  \item Active/decay patch ratio.
  \item Actual missing-force network call 수.
  \item Kernel launch 수와 p50/p95 dynamics latency.
  \item Patch toggle count와 activation duration.
  \item Activation acceleration/jerk.
\end{itemize}

\subsection*{공정한 비교}

두 방식으로 비교한다.

\begin{enumerate}[leftmargin=1.7em]
  \item 동일 active-patch budget에서 가장 낮은 dynamics error.
  \item 동일 p95 latency budget에서 가장 낮은 dynamics error.
\end{enumerate}

Gate classification score만 높고 wall-clock가 줄지 않으면 adaptive-compute claim은 성립하지 않는다.

\subsection*{결정 규칙}

\killgate Learned gate가 curvature/strain heuristic과 같은 비용에서 동등하면
learned gate를 main contribution에서 제외한다.
Always-on Local 대비 의미 있는 speedup 없이 error만 증가하면 adaptive claim을 철회한다.

\section{E6: Complement와 overlap conservation}

\subsection*{검증 claim}

C6: Local force/state를 Global complement로 제한하고
overlap을 보존적으로 조립해야 double counting과 seam이 줄어드는가?

\subsection*{Ablation}

\begin{enumerate}[leftmargin=1.7em]
  \item No complement.
  \item Runtime post-projection only.
  \item Offline complement basis only.
  \item Full-force target instead of missing-force target.
  \item Patch integrate-then-average.
  \item Partition-of-unity displacement only.
  \item Force assembly without torque constraint.
  \item Proposed complement + conservative force assembly.
\end{enumerate}

\subsection*{지표}

\begin{itemize}[leftmargin=1.5em]
  \item \(\|\Phi^\top M\Psi\|\)와 Global-span force leakage.
  \item Global modal amplitude inflation.
  \item Patch boundary position/velocity jump.
  \item Net force/torque error.
  \item Artificial center-of-mass drift/rotation.
  \item Energy jump at activation.
  \item Patch partition 변화에 대한 sensitivity.
\end{itemize}

\subsection*{필수 시각화}

Active core/halo map, overlap force arrows,
Global component와 Local complement component,
seam acceleration, net wrench curve를 한 sequence에 나란히 표시한다.

\section{E7: Local-to-Global feedback}

\subsection*{검증 claim}

C6의 feedback 부분:
Local correction이 visual overlay가 아니라 Global aeroelastic state를 실제로 바꾸는가?

\subsection*{비교군}

\begin{enumerate}[leftmargin=1.7em]
  \item No feedback: Local displacement는 rendering에만 추가.
  \item Structural cross-block feedback only.
  \item Corrected-aero delta only.
  \item Next-frame feedback.
  \item Same-frame one corrector.
  \item Same-frame multiple fixed-point iterations.
  \item Proposed minimal valid combination.
\end{enumerate}

\subsection*{지표}

\begin{itemize}[leftmargin=1.5em]
  \item Global modal coordinate와 tip trajectory 차이.
  \item Corrected aerodynamic generalized-force delta.
  \item Phase/amplitude와 teacher error.
  \item Extra latency와 iteration 수.
  \item Base-force double-count test.
\end{itemize}

\subsection*{결정 규칙}

\killgate Feedback 제거가 모든 benchmark에서 영향을 거의 주지 않으면
two-way coupling을 contribution으로 주장하지 않고
Local visual/kinematic corrector로 단순화한다.
Multiple corrector가 one-corrector보다 이득이 미미하면 한 번만 사용한다.

\section{E8: End-to-end static-3DGS comparison}

\subsection*{검증 claim}

C7: Static GS 입력에서 target mesh/full runtime continuum solver 없이
유용한 dynamics와 rendering Pareto를 제공하는가?

\subsection*{직접 또는 matched 비교 우선순위}

\begin{enumerate}[leftmargin=1.7em]
  \item Same-anchor sparse PD/XPBD.
  \item Gaussian Swaying-style surface aero + matched MPM.
  \item FreeForm-style mesh-free ROM.
  \item MeshGraphNets-style direct rollout.
  \item PGRD-style all-node residual.
  \item Proposed Global-only, always-on Local, oracle-gated, learned-gated.
\end{enumerate}

\subsection*{조건부 subset 비교}

\begin{itemize}[leftmargin=1.5em]
  \item DynamicTree: tree/leaf/branch-like subset.
  \item DiffWind: 공통 spatial wind forward-simulation subset.
  \item PhysSkin: 동일 reduced deformation/external load subset.
  \item PhysGaussian: 공통 force/material 설정 subset.
  \item GausSim/NGFF: 입력과 action/force를 공통화할 수 있을 때만.
\end{itemize}

\subsection*{Main table 항목}

\begin{itemize}[leftmargin=1.5em]
  \item Runtime target mesh 필요 여부.
  \item Runtime particle/grid continuum solver 필요 여부.
  \item Dynamics position/normal/frequency error.
  \item Rendering PSNR/SSIM/LPIPS와 temporal error.
  \item Setup time, dynamics p50/p95, transport/render time.
  \item Peak GPU memory.
  \item Unseen wind/object/material 결과.
\end{itemize}

\section{E9: Real static-GS transfer}

\subsection*{Capture}

최소 3개 real asset에 대해 canonical multi-view capture와 static 3DGS를 만든다.
가능하다면 fan/wind generator의 속도 profile,
고정점, object scale, material mass를 기록한다.
동적 ground truth는 별도 high-speed or multi-view video에서 얻는다.

\subsection*{정량 가능 항목}

\begin{itemize}[leftmargin=1.5em]
  \item Tracked marker/feature trajectory.
  \item Silhouette IoU.
  \item Tip amplitude와 dominant frequency.
  \item Wind-stop damping envelope.
  \item Novel-view temporal consistency.
\end{itemize}

정확한 force/material ground truth가 없으면
synthetic physics table과 real qualitative table을 분리한다.

\section{E10: Quality--latency와 scaling}

\subsection*{조절 변수}

\begin{itemize}[leftmargin=1.5em]
  \item Global rank \(r_g\): 예를 들어 8/16/32/64.
  \item Aerodynamic sample 수.
  \item Local basis rank.
  \item Patch core/halo 크기.
  \item 최대 active ratio.
  \item Gate threshold와 budget.
  \item Corrector iteration 수.
  \item Gaussian 수와 object batch 수.
\end{itemize}

\subsection*{보고할 curve}

\begin{itemize}[leftmargin=1.5em]
  \item Trajectory error 대 p95 dynamics latency.
  \item Spectrum error 대 p95 latency.
  \item Generalized-force/wrench error 대 aerodynamic sample 수.
  \item Active ratio 대 error와 latency.
  \item LPIPS 대 end-to-end latency.
  \item Object/Gaussian 수 대 memory와 throughput.
\end{itemize}

\subsection*{Latency breakdown}

\begin{enumerate}[leftmargin=1.7em]
  \item Surface reconstruction and wind query.
  \item Aerodynamic force and hyper-reduction.
  \item Global predictor.
  \item Gate.
  \item Missing-force network.
  \item Overlap/complement assembly.
  \item Local integrator.
  \item Global corrector.
  \item Gaussian transport.
  \item Rendering.
\end{enumerate}

Warm-up 후 GPU synchronization을 사용하고 p50/p95를 보고한다.
Batch 1과 multi-object batch, peak memory, kernel 수를 함께 기록한다.

\section{전체 ablation registry}

\begin{longtable}{L{0.20\linewidth}L{0.37\linewidth}L{0.12\linewidth}L{0.22\linewidth}}
\toprule
분류 & Ablation & Claim & 예상 실패 신호 \\
\midrule
\endfirsthead
\toprule
분류 & Ablation & Claim & 예상 실패 신호 \\
\midrule
\endhead
Topology & Oracle / proposed / kNN / FPS & C1 & Layer shortcut, mode error, rollout drift \\
\addlinespace
Topology & Oriented normal 제거 & C1,C3 & Front/back force error, normal flip \\
\addlinespace
Topology & Bending / stretching 구분 제거 & C1,C2 & Frequency / material control error \\
\addlinespace
Aero & Current 대신 rest normal/area & C3 & Large deformation force error \\
\addlinespace
Aero & Relative 대신 absolute wind & C3 & Phase/amplitude error \\
\addlinespace
Aero & Spatial 대신 object-average wind & C3 & Traveling gust failure \\
\addlinespace
Sampling & Full, random, FPS, nonconservative, proposed & C3,C7 & Wrench drift / latency tradeoff \\
\addlinespace
Global & Physics, hybrid, direct, recurrent & C2 & OOD / energy / stability difference \\
\addlinespace
Global & Mode count sweep & C2,C4,C7 & Local necessity and Pareto \\
\addlinespace
Local & None, full, reduced, hybrid & C4 & Flutter / fold recovery \\
\addlinespace
Local & Force residual 대 displacement residual & C4 & Wind-stop/phase stability \\
\addlinespace
Local & Temporal memory 제거 & C4 & Unsteady/history subset failure \\
\addlinespace
Gate & Never, always, random, heuristic, oracle, learned & C5 & Benefit regret and compute \\
\addlinespace
Gate & Binary current error 대 future benefit & C5 & Late activation and drift \\
\addlinespace
Gate & Hysteresis/cross-fade 제거 & C5,C6 & Chattering and jerk \\
\addlinespace
Coupling & Complement 제거 & C6 & Amplitude double counting \\
\addlinespace
Coupling & Conservative overlap 제거 & C6 & Seam, net torque drift \\
\addlinespace
Coupling & Structural reaction 제거 & C6 & Global response mismatch if cross block nonzero \\
\addlinespace
Coupling & Corrected aero 제거 & C6 & Local normal change ignored \\
\addlinespace
Coupling & Corrector 제거 & C6 & One-way feedback gap \\
\addlinespace
Transport & Center-only, rigid, affine & C7 & Covariance / orientation artifact \\
\bottomrule
\end{longtable}

\section{Generalization과 stress-test protocol}

\subsection{Generalization matrix}

\begin{longtable}{L{0.22\linewidth}L{0.31\linewidth}L{0.39\linewidth}}
\toprule
Split & 학습에 없는 조건 & 검증 항목 \\
\midrule
\endfirsthead
\toprule
Split & 학습에 없는 조건 & 검증 항목 \\
\midrule
\endhead
Direction & Unseen yaw/elevation & Directional aero and Global response \\
\addlinespace
Magnitude & Weak/strong extrapolation & Force scale, stability, gate recall \\
\addlinespace
Temporal & Unseen gust frequency/chirp & Resonance, phase, history \\
\addlinespace
Spatial & Moving/local gust & Hyper-reduction and active patch migration \\
\addlinespace
Material & Stiffness/damping/density & Physical editability \\
\addlinespace
Attachment & New support layout & Topology/operator/control transfer \\
\addlinespace
Geometry & New shape within category & Object generalization \\
\addlinespace
Category & Flag to curtain/paper/leaf & Cross-family limit \\
\addlinespace
Reconstruction & New GS fit/view count/seed & Target robustness \\
\addlinespace
Corruption & Floaters/holes/density change & Topology and binding robustness \\
\addlinespace
Combined OOD & New object+material+wind & Practical worst case \\
\bottomrule
\end{longtable}

\subsection{Stability stress tests}

\begin{longtable}{L{0.24\linewidth}L{0.33\linewidth}L{0.35\linewidth}}
\toprule
입력 & 확인할 현상 & 주요 지표 \\
\midrule
\endfirsthead
\toprule
입력 & 확인할 현상 & 주요 지표 \\
\midrule
\endhead
Zero-wind relaxation & Rest 복귀와 energy decay & Energy slope, residual displacement \\
\addlinespace
Abrupt start & 순간 가속/activation & Peak acceleration, jerk \\
\addlinespace
Abrupt stop & 복원과 damping & Envelope/frequency error \\
\addlinespace
Threshold-crossing wind & Gate flicker & Toggle count, jerk \\
\addlinespace
Constant strong wind & Steady drift/divergence & Energy growth, inversion \\
\addlinespace
Sinusoidal/chirp & Resonance/phase & Amplitude, dominant frequency \\
\addlinespace
Traveling gust & Active patch 이동 & Gate recall, spatial continuity \\
\addlinespace
30--60초 rollout & 누적 drift & Divergence rate, state norm \\
\addlinespace
\(\Delta t/2,2\Delta t\) & Integrator sensitivity & Trajectory/frequency change \\
\addlinespace
Strong OOD wind & Graceful failure & Covariance SPD, inversion, gate saturation \\
\bottomrule
\end{longtable}

\section{평가 지표}

\begin{longtable}{L{0.20\linewidth}L{0.31\linewidth}L{0.41\linewidth}}
\toprule
범주 & 지표 & 의미 \\
\midrule
\endfirsthead
\toprule
범주 & 지표 & 의미 \\
\midrule
\endhead
Geometry & Object-scale normalized position RMSE & 전체 deformation error \\
\addlinespace
State & Velocity/acceleration RMSE & Dynamic state error \\
\addlinespace
Surface & Normal angular error & Aero/render direction accuracy \\
\addlinespace
Surface & Strain/curvature error & Thin-shell response \\
\addlinespace
Dynamics & Tip amplitude/phase & Large response \\
\addlinespace
Dynamics & Dominant frequency/damping ratio & Oscillation and recovery \\
\addlinespace
Spectrum & PSD distance & Flutter/high-frequency detail \\
\addlinespace
Aero & Sample/generalized-force error & Aerodynamic operator accuracy \\
\addlinespace
Conservation & Net force/torque error & Cubature/overlap validity \\
\addlinespace
Energy & Work--energy residual/max growth & Nonphysical energy injection \\
\addlinespace
Coupling & Global-span leakage & Double-counted motion/force \\
\addlinespace
Patch & Boundary position/velocity jump & Seam artifact \\
\addlinespace
Gate & AUPRC/calibration/benefit regret & Selection quality \\
\addlinespace
Stability & Divergence/inversion rate & Long rollout robustness \\
\addlinespace
Gaussian & Invalid covariance/scale explosion & Render representation stability \\
\addlinespace
Rendering & PSNR/SSIM/LPIPS & Image fidelity \\
\addlinespace
Temporal render & Flow/feature warp error & Flicker and coherence \\
\addlinespace
Runtime & p50/p95, memory, kernel count & Actual interactive cost \\
\bottomrule
\end{longtable}

Rendered PSNR만으로 physics claim을 지지하지 않는다.
Main table에는 trajectory, frequency/damping, force/wrench, latency가 함께 있어야 한다.

\section{통계와 reporting}

\begin{itemize}[leftmargin=1.5em]
  \item Frame을 독립 sample로 취급하지 않는다.
  Sequence별 metric을 먼저 구하고 asset 단위로 집계한다.
  \item Asset bootstrap 95\% confidence interval을 보고한다.
  \item Learned baseline은 가능하면 최소 3개 seed를 사용한다.
  \item 모든 method에 동일 test trajectories와 initial conditions를 사용한다.
  \item Failure/divergence sequence를 평균에서 제외하지 않고 별도 비율과 capped metric으로 보고한다.
  \item 공식 구현이 아닌 style baseline은 implementation 차이를 supplementary에 기록한다.
  \item Runtime은 동일 hardware, precision, warm-up, synchronization 조건으로 측정한다.
\end{itemize}

\section{개발용 반증 기준과 go/no-go}

다음 수치는 논문 약속이 아니라 구현 중 방향을 유지할지 판단하는 임시 기준이다.

\begin{longtable}{L{0.31\linewidth}L{0.28\linewidth}L{0.33\linewidth}}
\toprule
질문 & 임시 통과 목표 & 실패 시 결정 \\
\midrule
\endfirsthead
\toprule
질문 & 임시 통과 목표 & 실패 시 결정 \\
\midrule
\endhead
Predicted topology가 충분한가? & Oracle dynamics error 대비 추가 악화 약 10--15\% 이내 & Oracle/simple topology scope 또는 target mesh baseline으로 축소 \\
\addlinespace
Local이 필요한가? & Global rank 증가보다 적은 비용으로 local spectrum/trajectory 개선 & Local contribution 제거/축소 \\
\addlinespace
Missing-force network가 필요한가? & Physics-only Local 대비 반복 가능한 개선 & Network 제거, pure adaptive physics \\
\addlinespace
Learned gate가 필요한가? & Heuristic보다 같은 budget에서 낮은 error & Learned gate contribution 제거 \\
\addlinespace
Adaptive Pareto가 있는가? & Always-on 대비 Local cost 약 2배 감소, dynamics error 증가 약 5\% 내외 & Adaptive claim 철회 \\
\addlinespace
End-to-end speedup이 있는가? & Always-on 대비 최소 약 1.25배 또는 same-quality 유의미한 개선 & Dispatcher/always-on/simple solver 전환 \\
\addlinespace
Global dynamics가 맞는가? & Tip frequency 약 5\%, amplitude 약 10\% 이내 & Basis/operator/aero 재설계 \\
\addlinespace
Long rollout이 안정적인가? & 4--10초 normalized position error 약 object scale 5\% 내외, divergence 없음 & Residual/gate rollout 재학습 또는 claim 축소 \\
\addlinespace
Conservative sampling이 필요한가? & Random/FPS보다 wrench/drift와 Pareto 개선 & Supporting detail로 하향 \\
\addlinespace
Feedback이 필요한가? & One-way 대비 measurable Global response/teacher error 개선 & Feedback claim 제거 \\
\bottomrule
\end{longtable}

\section{논문 결과표와 figure 계획}

\begin{enumerate}[leftmargin=1.7em]
  \item \textbf{Figure 1:} Static GS \(\rightarrow\) topology/operator package
  \(\rightarrow\) Global predictor \(\rightarrow\) active Local patches
  \(\rightarrow\) corrected Gaussians.
  \item \textbf{Table 1:} Oracle-topology matched solver accuracy/latency.
  \item \textbf{Table 2:} Static-GS end-to-end comparison과 input privilege.
  \item \textbf{Table 3:} Aero full/rest/random/FPS/RSH/proposed hyper-reduction.
  \item \textbf{Table 4:} Local physics/missing-force/gate/coupling ablations.
  \item \textbf{Table 5:} OOD와 long-rollout stability.
  \item \textbf{Figure 2:} Current-state aero와 rest-state WPB-style force 차이.
  \item \textbf{Figure 3:} Teacher/Global-only/physics-only Local/hybrid Local의 time response와 PSD.
  \item \textbf{Figure 4:} Oracle benefit map, learned gate, actual active set.
  \item \textbf{Figure 5:} Complement/overlap conservation 제거 시 seam, drift, torque error.
  \item \textbf{Figure 6:} Quality--latency Pareto와 active-ratio break-even.
  \item \textbf{Figure 7:} Oracle topology 대 predicted topology downstream motion.
  \item \textbf{Figure 8:} Real static-GS wind/material/attachment editing.
\end{enumerate}

\section{주요 위험과 실패 모드}

\subsection{Target-mesh-free 표현이 사실상 mesh가 되는 위험}

Oriented anchor graph가 watertight triangle connectivity와 full shell elements를 요구하면
reviewer는 target mesh를 다른 이름으로 부른 것이라고 볼 수 있다.
Runtime package는 sparse, non-watertight, locally oriented material graph로 정의하고,
mesh extraction baseline과 representation/storage/solver 차이를 명시해야 한다.
반대로 stable area-normal과 bending operator를 위해 실제 triangles가 필수라면
target-mesh-free claim을 약화하고 lightweight inferred proxy로 정직하게 표현한다.

\subsection{Topology distillation이 dynamics bottleneck이 되는 위험}

Close disconnected layers, thin tips, missing backside, floaters는
Euclidean graph와 orientation 추론을 쉽게 깨뜨린다.
Topology confidence가 낮을 때에는

\begin{itemize}[leftmargin=1.5em]
  \item uncertain edges 제거,
  \item user hint 또는 category prior 요청,
  \item conservative stiffness,
  \item target mesh extraction fallback
\end{itemize}

을 고려한다.
Fallback을 사용한 sequence는 main target-mesh-free 결과와 구분한다.

\subsection{Metric scale와 material identifiability}

일반 SfM/3DGS는 절대 scale이 없고,
wind speed, density, thickness, Young's modulus는 서로 보상될 수 있다.
User-supplied scale marker, known object dimension,
또는 nondimensional response preset 중 하나를 runtime contract에 넣는다.
실측 물성 추정과 artist-friendly control은 별도 평가한다.

\subsection{Global basis만으로 충분할 위험}

Global rank를 조금 늘리면 local flutter까지 충분히 재현될 수 있다.
이 경우 Local/gate/coupling 복잡성은 정당화되지 않는다.
E4에서 mode-count Pareto를 가장 먼저 수행하고,
Local이 같은 latency에서 분명한 high-frequency 이득을 보일 때만 유지한다.

\subsection{Local physics-only로 충분할 위험}

Teacher-S가 runtime과 같은 aerodynamic law와 structural law를 쓰면
learned missing force가 학습할 defect가 작을 수 있다.
Network가 이득을 주지 않으면 pure adaptive local physics가 더 깨끗한 논문이 된다.
Unsteady force를 인위적으로 만들기 위해 teacher와 runtime physics를 불필요하게 다르게 하지 않는다.

\subsection{Gate 비용이 절감량을 먹는 위험}

모든 patch에 heavy GNN gate를 실행하면 Local skip 이득이 사라진다.
Gate는 pooled Global state와 cheap local invariants를 사용하는 작은 model이어야 한다.
Gate latency, feature construction, gather/scatter를 별도 profile한다.

\subsection{Active basis 변화와 state remapping}

Patch가 활성화될 때 Local basis dimension이 바뀐다.
기존 \(z,\dot z\)를 새 basis로 energy-consistent하게 transfer해야 한다.

\begin{equation}
  z_{\mathrm{new}}
  =
  \argmin_z
  \left\|
  \Psi_{\mathrm{new}}z
  -
  \Psi_{\mathrm{old}}z_{\mathrm{old}}
  \right\|_M^2.
\end{equation}

New coordinate는 0 또는 boundary-consistent initialization으로 시작한다.
State transfer가 energy jump를 만들면 fixed superset Local basis와 sparse masks로 단순화한다.

\subsection{Missing-force inverse dynamics noise}

Position trajectory를 두 번 미분한 acceleration은 noisy하다.
가능하면 teacher force breakdown을 직접 저장하고,
그렇지 않으면 smoothing, weak-form/integral loss,
differentiable integrator를 통한 multi-step state supervision을 같이 쓴다.

\subsection{Overlap 보존과 complement constraint 충돌}

Wrench constraint와 Global complement를 순차 projection하면
한 constraint가 다른 constraint를 다시 깨뜨릴 수 있다.
필요하면 식~\eqref{eq:conservative-assembly}에서
wrench와 \(\Phi^\top f=0\)를 하나의 KKT system에 동시에 넣는다.
Constraint feasibility와 rank를 patch 구성마다 검사한다.

\subsection{Feedback이 수치적으로 불안정할 위험}

Local correction 후 aero를 다시 계산하면 follower-force feedback이 커질 수 있다.
One-corrector, under-relaxation, force-delta clamp,
smaller timestep을 stress test한다.
Full fixed-point convergence를 보장하지 않는다면
method를 predictor--one-corrector approximation으로 명시한다.

\subsection{Teacher와 real wind의 domain gap}

Quasi-steady teacher만으로 real cloth flutter의 phase와 spectrum을 맞추기 어렵다.
Real sequence는 system identification 또는 limited Teacher-FSI fine-tuning이 필요할 수 있다.
Synthetic physics accuracy와 real visual plausibility를 같은 claim으로 섞지 않는다.

\subsection{Gaussian transport와 physical surface 불일치}

Visual Gaussian support가 실제 material surface를 벗어나거나
두 겹 surface를 동시에 표현하면 동일 \(F\) transport가 artifact를 만든다.
Surface reliability, layer separation, covariance clamp,
outlier Gaussian fallback을 둔다.

\section{P0 implementation kill gates}

\begin{longtable}{L{0.06\linewidth}L{0.45\linewidth}L{0.41\linewidth}}
\toprule
\# & 중단 조건 & 필요한 조치 \\
\midrule
\endfirsthead
\toprule
\# & 중단 조건 & 필요한 조치 \\
\midrule
\endhead
1 & \(\|\Phi^\top M\Psi\|\)가 tolerance를 넘는다. & Basis construction/complement를 고치기 전 Local training 금지. \\
\addlinespace
2 & \(M\succ0,K,C\succeq0\) 또는 block symmetry가 깨진다. & Runtime solve 금지; operator construction 수정. \\
\addlinespace
3 & Gravity, attachment, aero, restoring, residual 중 owner가 둘 이상이다. & Force ledger를 다시 만들고 residual label 재생성. \\
\addlinespace
4 & Gate-off/zero-wind에서 rest로 돌아오지 않거나 Local energy가 증가한다. & Activation equation, damping, state decay 수정. \\
\addlinespace
5 & Patch별 integrate-then-average가 구현되어 있다. & Force assembly 후 단일 Local update로 변경. \\
\addlinespace
6 & 조립 후 Global leakage, overlap seam, internal net force/torque가 tolerance를 넘는다. & Joint constraint solve 또는 basis/patch 수정. \\
\addlinespace
7 & Global corrector가 corrected-minus-predictor delta 외 base force를 다시 사용한다. & Double-counting 제거. \\
\addlinespace
8 & Inactive 전환 시 \(z,\dot z\)를 즉시 freeze/reset한다. & Decay state와 energy threshold 추가. \\
\addlinespace
9 & Rigid/affine transport, covariance SPD, aero/render normal consistency가 실패한다. & Gaussian 결과 생성 금지; transport 수정. \\
\addlinespace
10 & Oracle Local이 Global-only를 개선하지 못한다. & Local branch 제거 또는 Global basis 재설계. \\
\addlinespace
11 & Always-on missing-force가 physics-only Local을 개선하지 못한다. & Network 제거; gate 학습 중단. \\
\addlinespace
12 & Learned gate가 end-to-end Pareto를 개선하지 못한다. & Adaptive claim 제거; always-on/pure physics 채택. \\
\bottomrule
\end{longtable}

\section{MVP와 구현 순서}

\subsection{MVP-A: Oracle-topology Global strip}

\begin{itemize}[leftmargin=1.5em]
  \item 한쪽 고정 cantilever strip/flag.
  \item Oracle topology, mass, stiffness, attachment.
  \item Current-surface analytic aero with all samples.
  \item \(r_g=8\)--32 Global modes.
  \item Global physics, Gaussian affine transport.
\end{itemize}

통과 조건은 start/stop, sinusoidal/chirp의 frequency/damping과
same-anchor sparse solver 대비 latency/accuracy다.

\subsection{MVP-B: Conservative aero reduction}

\begin{itemize}[leftmargin=1.5em]
  \item Full samples, random/FPS, nonconservative learned, proposed cubature.
  \item Held-out deformation과 traveling gust.
  \item Generalized force, net wrench, GPU latency.
\end{itemize}

이 단계에서 direct analytic과 RSH/LUT/MLP 비교도 끝낸다.

\subsection{MVP-C: Oracle active Local}

\begin{itemize}[leftmargin=1.5em]
  \item Oracle Local basis와 oracle active patch.
  \item Physics-only Local, full Local upper bound.
  \item Missing-force network는 initially always-on.
  \item Complement/overlap/feedback ablation.
\end{itemize}

Oracle Local이 Global mode 증가보다 나은 Pareto를 만들지 못하면 Local 방향을 중단한다.

\subsection{MVP-D: Learned gate}

\begin{itemize}[leftmargin=1.5em]
  \item Counterfactual benefit label.
  \item Lightweight gate, hard runtime skip.
  \item Active/decay/dormant state.
  \item Always/heuristic/oracle/learned 비교.
\end{itemize}

\subsection{MVP-E: Predicted target package}

\begin{itemize}[leftmargin=1.5em]
  \item Independent static GS variants.
  \item Oriented topology, area/mass, attachment, operator distillation.
  \item Oracle/predicted/kNN/mesh extraction comparison.
  \item Real static-GS qualitative transfer.
\end{itemize}

\subsection{첫 논문 권장 최소 범위}

\begin{itemize}[leftmargin=1.5em]
  \item Attached single-layer flag/cloth strip/paper/leaf.
  \item Prescribed spatial wind.
  \item Teacher-S 중심; Teacher-FSI는 작은 optional subset.
  \item Target-runtime sparse oriented anchor graph, no full mesh/MPM.
  \item Global physics + Local missing force + gate.
  \item Affine GS transport와 standard rendering.
\end{itemize}

Full tree, self-contact, tearing, free flight, appearance residual,
arbitrary cross-family material generalization은 후속으로 둔다.

\section{논문 framing}

\subsection{권장 제목}

\begin{quote}
\textbf{Topology-Distilled, Error-Triggered Global--Local Wind Dynamics
for Static 3D Gaussian Thin Surfaces}
\end{quote}

파일명은 제목의 핵심을 반영해

\begin{center}
\path{3dgs_topology_distilled_error_triggered_wind_dynamics_2026-08-13.tex}
\end{center}

로 정한다.

\subsection{대안 제목}

\begin{enumerate}[leftmargin=1.7em]
  \item \emph{Target-Mesh-Free Aeroelastic Animation of Static Gaussian Thin Surfaces
  with Error-Gated Complementary Dynamics}.
  \item \emph{Physics-Global, Learned-Local Wind Response for Static 3D Gaussian Thin Surfaces}.
  \item \emph{Conservative Reduced Wind Dynamics with Selective Local Refinement for Static 3DGS}.
\end{enumerate}

\subsection{안전한 novelty 문장}

\begin{quote}
To our knowledge, no prior wind-driven static-3DGS method couples
current-geometry Gaussian-surface loads to a topology-distilled object-level
reduced solver while selectively evaluating a future-error-benefit-gated
missing-force correction in the Global modal complement and feeding the
corrected aerodynamic state back into the Global dynamics.
\end{quote}

이 문장은 동일한 전체 input/output/coupling chain을 조사 범위에서 찾지 못했다는 의미다.
개별 modal equation, drag law, residual, adaptive subspace,
Gaussian transport가 최초라는 뜻이 아니다.

\subsection{권장 contribution 문장}

\begin{enumerate}[leftmargin=1.7em]
  \item A topology-distilled, target-runtime-mesh-free reduced physical representation
  for independently reconstructed static Gaussian thin surfaces.
  \item A current-state spatial aerodynamic hyper-reduction that preserves
  generalized loads and resultant force/torque.
  \item A future-benefit-gated complementary Local corrector that predicts
  only missing force, integrates it through Local physics, and closes the loop
  through conservative assembly and corrected-aero feedback.
\end{enumerate}

\subsection{피해야 할 표현}

\begin{itemize}[leftmargin=1.5em]
  \item First global--local Gaussian dynamics.
  \item First modal wind simulation for 3DGS.
  \item First Gaussian aerodynamic force.
  \item First adaptive local physics.
  \item First physics + neural residual 3DGS.
  \item Full fluid--structure interaction.
\end{itemize}

\section{최종 method 요약}

Training에서는 source thin-surface meshes와 high-resolution structural simulations를
privileged teacher로 사용한다.
Multi-view rendering에서 독립적으로 복원한 static 3DGS를 입력으로,
oriented material anchor graph, conservative area/mass,
attachment, bending/stretching operators, Global modes,
Local complement modes, aerodynamic cubature,
Gaussian binding을 distill한다.

Target runtime에서는 이전 Global/Local state에서 현재 anchor 위치, 속도,
normal, area를 복원하고 prescribed spatial wind와의 relative velocity로
quasi-steady aerodynamic sample force를 계산한다.
소수 sample의 force를 net force/torque와 object generalized force가 보존되도록 축약하고,
항상 켜진 Global mass--stiffness--damping system을 적분하여 큰 sway와 bending을 예측한다.

작은 gate network는 Global-only predictor, local strain/curvature,
wind, uncertainty, history에서 각 patch의 future correction benefit을 예측한다.
선택된 patch에서만 missing-force network를 실행하고,
같은 active/halo 영역의 analytic Local base aero와 함께 세 channel을 만든다.
Network는 위치가 아니라 base physics가 놓친 local force를 출력한다.
세 channel은 합치기 전에 각각 보존적으로 조립되고 Global span을 제거한다.
그 뒤 하나의 Local complement system이 관성, restoring, damping과 함께
channel별 generalized force의 합을 적분한다.

Local correction은 렌더링에만 더해지지 않는다.
수정된 surface normal, area, velocity에서 aerodynamic force를 다시 계산하고,
predictor 대비 aerodynamic delta와
updated coupling에서 predictor-consumed coupling을 뺀 structural cross delta만
작은 Global corrector에 되돌린다.
최종 \(q,z\)에서 anchor deformation을 한 번만 복원하고,
affine MLS로 Gaussian mean과 covariance를 운반하여 standard 3DGS로 렌더링한다.
Inactive patch의 expensive network는 실제 skip하지만,
이미 남은 Local energy는 물리적인 restoring/damping으로 decay한 뒤 종료한다.

\section{최종 연구 가치 판정}

이 방향의 강점은 Global/Local이라는 고전적 계층화 자체가 아니다.
각 레벨이 무엇을 입력으로 받고 무엇을 출력하는지,
어떤 force를 known physics가 소유하고 어떤 defect만 network가 학습하는지,
patch overlap과 Global complement를 어떻게 보존적으로 묶는지,
Local geometry 변화가 어떻게 aerodynamic load를 통해 다시 Global에 들어가는지를
wind-driven static 3DGS라는 target domain에 맞게 하나의 operator chain으로 정의한다는 점이다.

현재 가장 강한 novelty 후보는 Local gate 하나가 아니라 다음 결합이다.

\begin{equation}
\boxed{
\begin{aligned}
&\text{topology-distilled static-GS thin-surface ROM}\\
+\;&\text{current-state conservative aerodynamic reduction}\\
+\;&\text{future-benefit-gated complementary missing-force physics}\\
+\;&\text{closed-loop corrected-aero feedback}.
\end{aligned}
}
\end{equation}

그러나 exact chain이 선행에 없다는 사실만으로 충분하지 않다.
Gaussian Swaying, Wind Projection Basis, FreeForm/DynamicTree,
PGRD, Subspace Condensation/Trading Spaces의 단순 조합처럼 보이지 않으려면
E1--E7에서 topology operator, conservative aero,
benefit gate, complement/feedback 각각의 비자명한 필요성과
quality--latency Pareto를 반증 중심으로 보여야 한다.

\begingroup
\sloppy
\bibliographystyle{plain}
\bibliography{refs_topology_distilled_error_triggered_wind_dynamics}
\endgroup

\end{document}



\begin{itemize}[leftmargin=1.5em]
  \item \(G^0=\{(\mu_j^0,\Sigma_j^0,\alpha_j,c_j)\}_{j=1}^{N_G}\):
  독립적으로 재구성된 canonical static 3DGS.
  \item \(W(\mathbf x,t)\):
  사용자가 주거나 외부 시스템이 제공하는 prescribed spatial wind field.
  \item \(\vartheta\):
  density, stretching stiffness, bending stiffness, damping, aerodynamic coefficients를 포함하는 material controls.
  \item \(\mathcal C\):
  hard/soft attachment, support, handle와 boundary controls.
  \item \(s_{\mathrm{metric}}\):
  SfM/3DGS scale ambiguity를 해소하기 위한 metric scale 또는 명시적인 nondimensional preset.
\end{itemize}

Target mesh, target FEM/MPM particles, source camera trajectory는 runtime 필수 입력이 아니다.
Static 3DGS에서 만든 sparse material scaffold는 사용하지만,
이는 target에서 재구성한 watertight simulation mesh가 아니라
학습된 oriented anchor graph와 local surface support의 묶음이다.

\subsection{Runtime 출력}

각 frame의 출력은 다음이다.

\begin{equation}
  \mathcal O_t
  =
  \left(
    G^t,\;
    q^t,\dot q^t,\;
    z^t,\dot z^t,\;
    \mathcal A_t,\;
    \mathcal D_t
  \right),
\end{equation}

여기서 \(q\)는 object-wide Global state,
\(z\)는 Global complement에 놓인 Local state,
\(\mathcal A_t\)는 active/sustaining patch set,
\(\mathcal D_t\)는 force, gate, conservation, latency를 기록하는 diagnostics다.
최종 Gaussian \(G^t\)는 standard 3DGS renderer에 바로 전달된다.

\subsection{Training-only privileged input}

Training asset에서는 다음 정보를 사용할 수 있다.

\begin{equation}
  \mathcal I_{\mathrm{train}}
  =
  \left(
    \mathcal M,\;
    \mathcal T^{\mathrm{shell/MPM}},\;
    \mathcal T^{\mathrm{fluid}}\;(\text{optional}),\;
    G^0_{\mathrm{independent}}
  \right).
\end{equation}

\(\mathcal M\)은 source mesh와 material topology,
\(\mathcal T^{\mathrm{shell/MPM}}\)은 high-resolution structural teacher trajectory,
\(\mathcal T^{\mathrm{fluid}}\)은 선택적인 CFD/LBM 또는 실측 unsteady-aero teacher다.
Training 3DGS는 source mesh vertices에 Gaussian을 그대로 부착하여 만들지 않고,
rendered multi-view images에서 독립적으로 reconstruct한 variant를 포함해야 한다.
그래야 target의 reconstruction noise, density variation, floaters, backside missing 문제를 실제로 학습한다.
